% Faithful initial reconstruction of Clause 12 from source PDF pages 303--346
% Native, searchable reconstructions of IEEE Std 1516.1-2010 Clause 12
% Tables 18--28.  Each command is deliberately page-local so the printed
% source pagination can remain fixed while the table text stays searchable.

\newcommand{\cxiitabletitle}[1]{%
    \par\vspace{2pt}%
    \begin{center}\sffamily\bfseries\fontsize{9.5}{10.9}\selectfont #1\end{center}%
    \vspace{3pt}%
}

\newenvironment{cxiitabletypeset}[2][8.0]{%
    \begingroup
    \fontsize{#1}{9.15}\selectfont
    \setlength{\tabcolsep}{5.0pt}%
    \renewcommand{\arraystretch}{#2}%
    \setlength{\arrayrulewidth}{0.45pt}%
    \sloppy
    \hbadness=10001\relax
}{%
    \par\endgroup
}

\newcommand{\cxiitableeighteen}{%
    \cxiitabletitle{Table 18---Standardized time types}%
    \begin{cxiitabletypeset}[8.20]{1.72}
    \makebox[\linewidth][c]{%
        \begin{tabular}{|Y{116pt}|Y{116pt}|Y{140pt}|}
        \hline
        \Centering\arraybackslash\bfseries Property &
        \Centering\arraybackslash\bfseries 64-bit integer time &
        \Centering\arraybackslash\bfseries 64-bit float time \\
        \hline
        Symbolic name & ``HLAinteger64Time'' & ``HLAfloat64Time'' \\
        \hline
        Representation & 64-bit integer & 64-bit float \\
        \hline
        Encoded format & HLAinteger64BE & HLAfloat64BE \\
        \hline
        Initial value & 0 & 0.0 \\
        \hline
        Final value & $2^{63}-1$ & 0x1.FFFFFFFFFFFFFP+1023 \\
        \hline
        Zero value & 0 & 0.0 \\
        \hline
        Epsilon value & 1 & 0x0.0000000000001P-1022 \\
        \hline
        \end{tabular}%
    }
    \end{cxiitabletypeset}%
}

\newcommand{\cxiitablenineteenandtwenty}{%
    \cxiitabletitle{Table 19---Noncomplex Java encoding helpers}%
    \begin{cxiitabletypeset}[8.00]{1.34}
    \makebox[\linewidth][c]{%
        \begin{tabular}{|Y{132pt}|Y{132pt}|Y{132pt}|}
        \hline
        \Centering\arraybackslash\bfseries HLA data\newline representation &
        \Centering\arraybackslash\bfseries Encoding helper class &
        \Centering\arraybackslash\bfseries Java type \\
        \hline
        HLAoctet & HLAoctet & byte \\ \hline
        HLAoctetPairBE & HLAoctetPairBE & short \\ \hline
        HLAoctetPairLE & HLAoctetPairLE & short \\ \hline
        HLAinteger16BE & HLAinteger16BE & short \\ \hline
        HLAinteger16LE & HLAinteger16LE & short \\ \hline
        HLAinteger32BE & HLAinteger32BE & int \\ \hline
        HLAinteger32LE & HLAinteger32LE & int \\ \hline
        HLAinteger64BE & HLAinteger64BE & long \\ \hline
        HLAinteger64LE & HLAinteger64LE & long \\ \hline
        HLAfloat32BE & HLAfloat32BE & float \\ \hline
        HLAfloat32LE & HLAfloat32LE & float \\ \hline
        HLAfloat64BE & HLAfloat64BE & double \\ \hline
        HLAfloat64LE & HLAfloat64LE & double \\ \hline
        HLAASCIIchar & HLAASCIIchar & byte \\ \hline
        HLAunicodeChar & HLAunicodeChar & short \\ \hline
        HLAbyte & HLAbyte & byte \\ \hline
        HLAhandle & HLAhandle & byte[] \\ \hline
        HLAlogicalTime & HLAlogicalTime & byte[] \\ \hline
        HLAlogicalTimeInterval & HLAbyte & byte[] \\ \hline
        HLAASCIIString & HLAASCIIString & String \\ \hline
        HLAunicodeString & HLAunicodeString & String \\ \hline
        HLAopaqueData & HLAopaqueData & byte[] \\ \hline
        \end{tabular}%
    }
    \end{cxiitabletypeset}%
    \vspace{8pt}%
    \cxiitabletitle{Table 20---Complex Java encoding helpers}%
    \begin{cxiitabletypeset}[8.00]{1.40}
    \makebox[\linewidth][c]{%
        \begin{tabular}{|Y{132pt}|Y{132pt}|}
        \hline
        \Centering\arraybackslash\bfseries HLA data\newline representation &
        \Centering\arraybackslash\bfseries Encoding helper class \\
        \hline
        HLAfixedRecord & HLAfixedRecord \\ \hline
        HLAvariantRecord & HLAvariantRecord \\ \hline
        HLAfixedArray & HLAfixedArray \\ \hline
        HLAvariableArray & HLAvariableArray \\ \hline
        \end{tabular}%
    }
    \end{cxiitabletypeset}%
}

\newcommand{\cxiitabletwentyone}{%
    \cxiitabletitle{Table 21---Noncomplex C++ encoding helpers}%
    \begin{cxiitabletypeset}[8.00]{1.34}
    \makebox[\linewidth][c]{%
        \begin{tabular}{|Y{132pt}|Y{132pt}|Y{132pt}|}
        \hline
        \Centering\arraybackslash\bfseries HLA data\newline representation &
        \Centering\arraybackslash\bfseries Encoding helper class &
        \Centering\arraybackslash\bfseries C++ type/macro \\
        \hline
        HLAoctet & HLAoctet & Octet \\ \hline
        HLAoctetPairBE & HLAoctetPairBE & OctetPair \\ \hline
        HLAoctetPairLE & HLAoctetPairLE & OctetPair \\ \hline
        HLAinteger16BE & HLAinteger16BE & Integer16 \\ \hline
        HLAinteger16LE & HLAinteger16LE & Integer16 \\ \hline
        HLAinteger32BE & HLAinteger32BE & Integer32 \\ \hline
        HLAinteger32LE & HLAinteger32LE & Integer32 \\ \hline
        HLAinteger64BE & HLAinteger64BE & Integer64 \\ \hline
        HLAinteger64LE & HLAinteger64LE & Integer64 \\ \hline
        HLAfloat32BE & HLAfloat32BE & float \\ \hline
        HLAfloat32LE & HLAfloat32LE & float \\ \hline
        HLAfloat64BE & HLAfloat64BE & double \\ \hline
        HLAfloat64LE & HLAfloat64LE & double \\ \hline
        HLAASCIIchar & HLAASCIIchar & char \\ \hline
        HLAunicodeChar & HLAunicodeChar & wchar\_t \\ \hline
        HLAbyte & HLAbyte & char \\ \hline
        HLAhandle & HLAhandle & char \\ \hline
        HLAlogicalTime & HLAlogicalTime & char \\ \hline
        HLAlogicalTimeInterval & HLAlogicalTimeInterval & char \\ \hline
        HLAASCIIString & HLAASCIIString & std::string \\ \hline
        HLAunicodeString & HLAunicodeString & std::wstring \\ \hline
        HLAopaqueData & HLAopaqueData & HLAopaqueData \\ \hline
        \end{tabular}%
    }
    \end{cxiitabletypeset}%
}

\newcommand{\cxiitabletwentytwo}{%
    \cxiitabletitle{Table 22---Complex C++ encoding helpers}%
    \begin{cxiitabletypeset}[8.10]{1.45}
    \makebox[\linewidth][c]{%
        \begin{tabular}{|Y{132pt}|Y{132pt}|}
        \hline
        \Centering\arraybackslash\bfseries HLA data\newline representation &
        \Centering\arraybackslash\bfseries Encoding helper class \\
        \hline
        HLAfixedRecord & HLAfixedRecord \\ \hline
        HLAvariantRecord & HLAvariantRecord \\ \hline
        HLAfixedArray & HLAfixedArray \\ \hline
        HLAvariableArray & HLAvariableArray \\ \hline
        \end{tabular}%
    }
    \end{cxiitabletypeset}%
}

\newcommand{\cxiitabletwentythree}{%
    \cxiitabletitle{Table 23---C++ API provided time implementations}%
    \begin{cxiitabletypeset}[7.80]{1.45}
    \makebox[\linewidth][c]{%
        \begin{tabular}{|Y{142pt}|Y{162pt}|Y{162pt}|}
        \hline
        \Centering\arraybackslash\bfseries Property &
        \Centering\arraybackslash\bfseries 64-bit integer time &
        \Centering\arraybackslash\bfseries 64-bit float time \\
        \hline
        Symbolic name & ``HLAinteger64Time'' & ``HLAfloat64Time'' \\ \hline
        C++ time class name & rti1516e::HLAinteger64Time & rti1516e::HLAfloat64Time \\ \hline
        C++ interval class name & rti1516e::HLAinteger64Interval & rti1516e::HLAfloat64Interval \\ \hline
        C++ time factory class name & rti1516e::HLAinteger64TimeFactory & rti1516e::HLAfloat64TimeFactory \\ \hline
        \end{tabular}%
    }
    \end{cxiitabletypeset}%
}

\newcommand{\cxiitabletwentyfour}{%
    \cxiitabletitle{Table 24---WSDL referenced standards}%
    \begin{cxiitabletypeset}[8.00]{1.45}
    \makebox[\linewidth][c]{%
        \begin{tabular}{|Y{154pt}|Y{292pt}|}
        \hline
        \Centering\arraybackslash\bfseries Specification & \Centering\arraybackslash\bfseries Location \\
        \hline
        WSDL 1.1 & {\urlstyle{same}\nolinkurl{http://www.w3.org/TR/wsdl.html}} \\ \hline
        SOAP 1.1 & {\urlstyle{same}\nolinkurl{http://www.w3.org/TR/2000/NOTE-SOAP-20000508/}} \\ \hline
        XML 1.0 & {\urlstyle{same}\nolinkurl{http://www.w3.org/TR/REC-xml/}} \\ \hline
        HTTP 1.1 & {\urlstyle{same}\nolinkurl{http://www.ietf.org/rfc/rfc2616}} \\ \hline
        XML Schema Part 1: Structures & {\urlstyle{same}\nolinkurl{http://www.w3.org/TR/xmlschema-1/}} \\ \hline
        XML Schema Part 2: Datatypes & {\urlstyle{same}\nolinkurl{http://www.w3.org/TR/xmlschema-2/}} \\ \hline
        WS-I Basic Profile 1.1 & {\urlstyle{same}\nolinkurl{http://www.ws-i.org/Profiles/BasicProfile-1.1-2004-08-24.html}} \\ \hline
        \end{tabular}%
    }
    \end{cxiitabletypeset}%
}

\newcommand{\cxiiwsseq}[1]{\textless{}xsd:sequence\textgreater{} #1}
\newcommand{\cxiiwsseqtight}[1]{\textless{}xsd:sequence\textgreater{}#1}
\newcommand{\cxiiwsenum}{\textless{}xsd:enumeration\textgreater{} of \textless{}xsd:string\textgreater{}}

\newcommand{\cxiitabletwentyfivea}{%
    \cxiitabletitle{Table 25---Parameter type representations}%
    \begin{cxiitabletypeset}[8.00]{1.92}
    \makebox[\linewidth][c]{%
        \begin{tabular}{|Y{220pt}|Y{300pt}|}
        \hline
        \Centering\arraybackslash\bfseries Parameter & \Centering\arraybackslash\bfseries Representation in Web Service API \\
        \hline
        userSuppliedTag & xsd:base64Binary \\ \hline
        timeStamp & xsd:base64Binary \\ \hline
        logicalTime & xsd:base64Binary \\ \hline
        lookahead & xsd:base64Binary \\ \hline
        messageRetractionDesignator & xsd:base64Binary \\ \hline
        regionDesignator & xsd:base64Binary \\ \hline
        attributeValue & xsd:base64Binary \\ \hline
        parameterValue & xsd:base64Binary \\ \hline
        encodedHandle & xsd:base64Binary \\ \hline
        passiveSubscriptionIndicator & xsd:boolean \\ \hline
        requestSuccessIndicator & xsd:boolean \\ \hline
        federationRestoreSuccessIndicator & xsd:boolean \\ \hline
        reservationSuccessIndicator & xsd:boolean \\ \hline
        registrationSuccessIndicator & xsd:boolean \\ \hline
        federationSaveSuccessIndicator & xsd:boolean \\ \hline
        instanceAttributeOwnershipIndicator & xsd:boolean \\ \hline
        timeIsValid & xsd:boolean \\ \hline
        \end{tabular}%
    }
    \end{cxiitabletypeset}%
}

\newcommand{\cxiitabletwentyfiveb}{%
    \cxiitabletitle{Table 25---Parameter type representations \textit{(continued)}}%
    \begin{cxiitabletypeset}[7.80]{1.90}
    \makebox[\linewidth][c]{%
        \begin{tabular}{|Y{220pt}|Y{300pt}|}
        \hline
        \Centering\arraybackslash\bfseries Parameter & \Centering\arraybackslash\bfseries Representation in Web Service API \\
        \hline
        dimensionBound & xsd:long \\ \hline
        rangeBound & xsd:long \\ \hline
        normalizedValue & xsd:long \\ \hline
        maximumCallbackCount & xsd:long \\ \hline
        pendingIndicator & xsd:long \\ \hline
        updateRate & xsd:double \\ \hline
        federateName & xsd:string \\ \hline
        federationExecutionName & xsd:string \\ \hline
        objectInstanceName & xsd:string \\ \hline
        objectClassName & xsd:string \\ \hline
        attributeName & xsd:string \\ \hline
        interactionClassName & xsd:string \\ \hline
        parameterName & xsd:string \\ \hline
        dimensionName & xsd:string \\ \hline
        transportationName & xsd:string \\ \hline
        orderName & xsd:string \\ \hline
        localSettingsDesignator & xsd:string \\ \hline
        federateDesignator & xsd:string \\ \hline
        objectClassDesignator & xsd:string \\ \hline
        interactionClassDesignator & xsd:string \\ \hline
        objectInstanceDesignator & xsd:string \\ \hline
        updateRateDesignator & xsd:string \\ \hline
        parameterDesignator & xsd:string \\ \hline
        attributeDesignator & xsd:string \\ \hline
        dimensionDesignator & xsd:string \\ \hline
        ownershipDesignator & xsd:string \\ \hline
        federateType & xsd:string \\ \hline
        synchronizationPointLabel & xsd:string \\ \hline
        federationSaveLabel & xsd:string \\ \hline
        faultDescription & xsd:string \\ \hline
        fomDocumentDesignator & xsd:anyURI \\ \hline
        AttributeDesignatorValuePair & \cxiiwsseq{of attributeDesignator and attributeValue} \\ \hline
        ParameterValuePair & \cxiiwsseq{of parameterDesignator and parameterValue} \\ \hline
        \end{tabular}%
    }
    \end{cxiitabletypeset}%
}

\newcommand{\cxiitabletwentyfivec}{%
    \cxiitabletitle{Table 25---Parameter type representations \textit{(continued)}}%
    \begin{cxiitabletypeset}[7.70]{1.88}
    \makebox[\linewidth][c]{%
        \begin{tabular}{|Y{220pt}|Y{300pt}|}
        \hline
        \Centering\arraybackslash\bfseries Parameter & \Centering\arraybackslash\bfseries Representation in Web Service API \\
        \hline
        AttributeDesignatorRegionDesignatorPair & \cxiiwsseq{of attributeDesignator and regionDesignator} \\ \hline
        FederateDesignatorSaveStatusPair & \cxiiwsseq{of federateDesignator and SaveStatus} \\ \hline
        FederateDesignatorRestoreStatusPair & \cxiiwsseq{of federateDesignator and RestoreStatus} \\ \hline
        TimeQueryReturn & \cxiiwsseq{of timeIsValid and logicalTime} \\ \hline
        RangeBounds & \cxiiwsseq{of two rangeBound} \\ \hline
        Callback & \textless{}xsd:choice\textgreater{} between one of the callbacks \\ \hline
        ResignAction & \cxiiwsenum \\ \hline
        OrderType & \cxiiwsenum \\ \hline
        TransportationType & \cxiiwsenum \\ \hline
        ServiceGroup & \cxiiwsenum \\ \hline
        SaveFailureReason & \cxiiwsenum \\ \hline
        SynchronizationPointFailureReason & \cxiiwsenum \\ \hline
        RestoreFailureReason & \cxiiwsenum \\ \hline
        SaveStatus & \cxiiwsenum \\ \hline
        RestoreStatus & \cxiiwsenum \\ \hline
        AttributeDesignatorSet & \cxiiwsseq{of attributeDesignator} \\ \hline
        AttributeDesignatorValuePairSet & \cxiiwsseq{of AttributeDesignatorValuePairs} \\ \hline
        ParameterDesignatorValuePairSet & \cxiiwsseqtight{of ParameterDesignatorValuePairs} \\ \hline
        DimensionDesignatorSet & \cxiiwsseqtight{of dimensionDesignators} \\ \hline
        RegionDesignatorSet & \cxiiwsseq{of regionDesignators} \\ \hline
        FederateDesignatorSet & \cxiiwsseq{of federateDesignators} \\ \hline
        AttributeDesignatorRegionDesignatorPairList & \cxiiwsseq{of AttributeDesignatorRegionDesignatorPairs} \\ \hline
        CallbackArray & \cxiiwsseq{of Callbacks} \\ \hline
        CallbackReturn & \cxiiwsseq{of CallbackArray and pendingIndicator} \\ \hline
        FederateDesignatorSaveStatusPairList & \cxiiwsseq{of FederateDesignatorSaveStatusPairs} \\ \hline
        FederateDesignatorRestoreStatusPairList & \cxiiwsseq{of FederateDesignatorRestoreStatusPairs} \\ \hline
        FederationExecutionNameSet & \cxiiwsseq{of federationExecutionNames} \\ \hline
        \end{tabular}%
    }
    \end{cxiitabletypeset}%
}

\newcommand{\cxiitabletwentysix}{%
    \cxiitabletitle{Table 26---WSDL standard time types}%
    \begin{cxiitabletypeset}[7.80]{1.45}
    \makebox[\linewidth][c]{%
        \begin{tabular}{|Y{118pt}|Y{154pt}|Y{154pt}|}
        \hline
        \Centering\arraybackslash\bfseries Property &
        \Centering\arraybackslash\bfseries 64-bit integer time\newline service name &
        \Centering\arraybackslash\bfseries 64-bit float time\newline service name \\
        \hline
        Symbolic name & HLAinteger64Time & HLAfloat64Time \\ \hline
        Encode time value & encodeHLAinteger64Time & encodeHLAfloat64Time \\ \hline
        Decode time value & decodeHLAinteger64Time & decodeHLAfloat64Time \\ \hline
        Encode interval value & encodeHLAinteger64Interval & encodeHLAfloat64Interval \\ \hline
        Decode interval value & decodeHLAinteger64Interval & decodeHLAfloat64Interval \\ \hline
        \end{tabular}%
    }
    \end{cxiitabletypeset}%
}

\newcommand{\cxiitabletwentysevena}{%
    \cxiitabletitle{Table 27---Handle name sources}%
    \begin{cxiitabletypeset}[7.80]{1.45}
    \makebox[\linewidth][c]{%
        \begin{tabular}{|Y{126pt}|Y{336pt}|}
        \hline
        \Centering\arraybackslash\bfseries Handle & \Centering\arraybackslash\bfseries Source for name \\
        \hline
        Object Class Handle & fully qualified name from FDD object class structure table \\ \hline
        Attribute Handle & fully qualified name from FDD attribute table \\ \hline
        Interaction Class Handle & fully qualified name from FDD interaction class structure table \\ \hline
        Parameter Handle & fully qualified name from FDD parameter table \\ \hline
        Dimension Handle & name from FDD dimension table \\ \hline
        Object Instance Handle & object instance name as returned from registerObjectInstance operation \\ \hline
        \end{tabular}%
    }
    \end{cxiitabletypeset}%
}

\newcommand{\cxiitabletwentysevenb}{%
    \cxiitabletitle{Table 27---Handle name sources \textit{(continued)}}%
    \begin{cxiitabletypeset}[7.85]{1.45}
    \makebox[\linewidth][c]{%
        \begin{tabular}{|Y{126pt}|Y{336pt}|}
        \hline
        \Centering\arraybackslash\bfseries Handle & \Centering\arraybackslash\bfseries Source for name \\
        \hline
        Federate Handle & federate name as provided to the joinFederationExecution operation \\ \hline
        Transportation Type & name from FDD transportation type table \\ \hline
        Order Type & name from FDD attribute table or FDD parameter table \\ \hline
        Update Rate & name from FDD update rate table \\ \hline
        \end{tabular}%
    }
    \end{cxiitabletypeset}%
}

\newcommand{\cxiitabletwentyeight}{%
    \cxiitabletitle{Table 28---Handle encoding and decoding services}%
    \begin{cxiitabletypeset}[7.80]{1.45}
    \makebox[\linewidth][c]{%
        \begin{tabular}{|Y{123pt}|Y{166pt}|Y{166pt}|}
        \hline
        \Centering\arraybackslash\bfseries Handle type &
        \Centering\arraybackslash\bfseries Encoding service &
        \Centering\arraybackslash\bfseries Decoding service \\
        \hline
        Federate Handle & \textit{Encode Federate Handle} & \textit{Decode Federate Handle} \\ \hline
        Object Class Handle & \textit{Encode Object Class Handle} & \textit{Decode Object Class Handle} \\ \hline
        Interaction Class Handle & \textit{Encode Interaction Class Handle} & \textit{Decode Interaction Class Handle} \\ \hline
        Object Instance Handle & \textit{Encode Object Instance Handle} & \textit{Decode Object Instance Handle} \\ \hline
        Parameter Handle & \textit{Encode Parameter Handle} & \textit{Decode Parameter Handle} \\ \hline
        Attribute Handle & \textit{Encode Attribute Handle} & \textit{Decode Attribute Handle} \\ \hline
        Dimension Handle & \textit{Encode Dimension Handle} & \textit{Decode Dimension Handle} \\ \hline
        \end{tabular}%
    }
    \end{cxiitabletypeset}%
}

% (standard pages 289--332). Source-page boundaries are intentionally preserved.

\clearpage

% --- Standard page 289 / source PDF page 303 ---
\begin{compactstdpage}
    \section{Programming language mappings}

    \subsection{Overview}

    The services described in this document to this point have been described in a programming-language-neutral way. This approach was taken because support is provided for multiple programming languages as described in the annexes to this document. This clause (i.e., Clause 12) provides the reader with a mapping from the abstract view from the body of this document (Clause 4 through Clause 11) to the programming-language-specific views in the APIs (Annex B through Annex D). It shall be transparent to the federation and participating federates which API or implementation language a specific federate uses.

    The mappings described in this clause are

    \stdliteraldashitem{Of designators from Clause 4 through Clause 10 to handles and names in the APIs}
    \stdliteraldashitem{From services in Clause 4 through Clause 10 to methods in the APIs}
    \stdliteraldashitem{From abstract datatypes in Clause 4 through Clause 10 to classes or datatypes in the APIs}

    \subsection{Designators}

    As explained in 1.4.3, designators are a generalization of both names and handles. In most cases, what is referred to as a designator in Clause 4 through Clause 10 has both a name and a handle. Which specific form of a designator is used in the APIs is dependent on the language and service in which it is used. Examine the APIs to see which form is needed, although in most cases, the handle is used. Clause 10 presents services that shall be used to map names to handles and vice versa.

    Names in most cases (e.g., object class names, attribute names, interaction class names) come from the FDD. As such, they are known ahead of time to all federates. Object instance names, however, are different. These names may be chosen by the registering federate, or they may be created dynamically by the RTI. As such, these names often are not predictable.

    Handles in all cases are generated by the RTI. In most cases, where the handles correspond to names from the FDD, these handles shall be generated in a repeatable manner. In other words, if two federation executions are created on the same version of the same RTI with the same FDD file, then the handles assigned to the same name shall be the same in each federation execution. This property facilitates saves and restores so that handles, rather than just the names, can be saved by federates. This property does not, however, imply that an RTI assign handles in a manner known a priori to federates.

    Handles shall be unique in a federation execution. Each federate shall know the same item by the same handle. In other words, the federate registering an object instance shall know it by the same handle as another federate that discovers the object instance.

    Often during a federation execution, it is useful for federates to send information identifying an item to another federate. For example, one federate may wish to send an interaction to another federate informing it that one particular object instance is of particular importance. To send information identifying an item to another federate, there are two options: send the name of the item, or send the handle of the item.

    To send a handle as an attribute value or as a parameter value, the handle shall be encoded prior to sending, and it shall be reconstituted by the receiver. The precise means of encoding and reconstituting a handle is described for each API in 12.11.5.4, 12.12.6.4, and 12.13.4. However, a handle encoded using any API (e.g., C++) shall be reconstitutable in all APIs. Also, a handle of a given type (e.g., an object instance handle) shall be reconstitutable only as that type.

\end{compactstdpage}

\clearpage

% --- Standard page 290 / source PDF page 304 ---
\begin{compactstdpage}
    Handles that are sent as part of MOM attribute value updates or as parameter values of MOM interactions are sent encoded and may be reconstituted in the manner described for each API.

    \subsection{Logical time, timestamps, and lookahead}

    In order to support customized implementations of logical time, timestamps, and lookahead, the C++ and Java APIs define a group of closely related abstract interfaces: LogicalTime, LogicalTimeInterval, and LogicalTimeFactory. These abstract interfaces have been designed to be implemented by the federate developer for use by an RTI implementation. This clause defines the general characteristics of the abstract interfaces, and each API annex presents API-specific details.

    The LogicalTime abstract interface is used to represent logical time and timestamp values. It has methods to compare the relative order of one instance with another and to compare an instance with the initial value (e.g., zero) and with the final value (e.g., infinity). This generality is provided to allow federate programmers to define arbitrary lower and upper bounds on their implementation of the LogicalTime class. Methods shall exist to add and subtract a value to and from of an instance of LogicalTime. However, it is noted that the argument to these methods is of type LogicalTimeInterval, and not of type LogicalTime.

    The LogicalTimeInterval abstract interface is similar to the LogicalTime abstract interface, but it is used to represent lookahead values. The chief differences in its interface are that there is no concept of a ``final'' or ``largest possible'' LogicalTimeInterval instance; instead, there are methods to test if a LogicalTimeInterval instance is zero or epsilon. Epsilon is the smallest nonzero LogicalTimeInterval that is possible with the particular implementation. This value is used by an RTI implementation to support time management services for federates with zero lookahead. The value of epsilon is intended to be constant. Methods shall exist to add and subtract a value to and from a LogicalTimeInterval instance.

    LogicalTime and LogicalTimeInterval abstract interfaces shall provide methods to obtain a string representation suitable for printing. The set of abstract interfaces for LogicalTime, LogicalTimeInterval, and LogicalTimeFactory shall also support some way to

    \stdliteraldashitem{Encode an instance of LogicalTime or LogicalTimeInterval into a compact, opaque representation suitable for network transmission into the provided variable array at the specified offset}
    \stdliteraldashitem{Create a LogicalTime or LogicalTimeInterval from the network representation in the provided variable array at the indicated offset.}
    The LogicalTimeFactory abstract interface shall have a method to create a LogicalTime instance whose value is the user-defined initial value and a method to create a LogicalTime instance whose value is the user-defined final value. The LogicalTimeFactory abstract interface shall have a method to create a LogicalTimeInterval instance whose value is zero and a method to create a LogicalTimeInterval instance whose value is the epsilon value.

    To allow the users to get the LogicalTimeFactory from the RTI, the C++ and Java APIs add a getTimeFactory() method to the RTI ambassador interface. Each API returns a value appropriate for the language.

    Implementers of the C++ and Java APIs shall provide, as part of their implementation, two implementations of the logical time abstract interfaces. One logical time abstract interface, the HLAInteger64Time, shall be based on the HLAInteger64BE type and the other, the HLAFloat64Time, shall be based on the HLAFloat64BE type. The following are intended to describe requirements on ALL time implementations.

\end{compactstdpage}

\clearpage

% --- Standard page 291 / source PDF page 305 ---
\begin{compactstdpage}
    Given time T and nonzero interval I, zero interval Iz, and epsilon interval Ie, all implementations of logical time shall ensure that the following statements are true:

    \begin{flushleft}
        \hspace*{1.2em}\begin{tabular}{@{}l@{}}
                           $T' = T + I \Rightarrow T < T'$               \\
                           $T' = T - I \Rightarrow T > T'$               \\
                           $T' = T + I_z \Rightarrow T = T'$             \\
                           $T' = T - I_z \Rightarrow T = T'$             \\
                           $T_1 \ne T_2 \Rightarrow |T_1 - T_2| \ge I_e$ \\
                           $T_1 = T_2 \Rightarrow T_2 - T_1 = I_z$       \\
                           $I' = I + I_z \Rightarrow I = I'$             \\
                           $I' = I - I_z \Rightarrow I = I'$             \\
                           $I_2 \ne I_1 \Rightarrow |I_2 - I_1| \ge I_e$ \\
                           $I_2 = I_1 \Rightarrow I_2 - I_1 = I_z$       \\
                           $I \ne I_z \Rightarrow I \ge I_e$             \\
                           $T > T' \Rightarrow T' < T$                   \\
        \end{tabular}
    \end{flushleft}

    \subsection{Standardized time types}

    An RTI shall support the standardized time representation shown in Table 18.

    \cxiitableeighteen

    When a federate execution is created, the time representation used is specified using the symbolic name above. All time representations shall provide methods through which a time value can be increased and decreased by epsilon while guaranteeing that the increased value differs from the original value. The epsilon value for HLAinteger64Time shall equal 1.

    \subsubsection{HLAInteger64Time}

    The implementation of HLAInteger64Time shall provide the following:

    \stdliteraldashitem{An implementation of the interface LogicalTime called HLAInteger64Time. This implementation shall use an HLAInteger64BE as its underlying representation of time. It shall use 0 as the value of ``initialTime'' and the largest positive value defined for an HLAInteger64BE as its final time. Methods that are defined to return or accept a LogicalTime reference shall actually return or expect a reference to HLAInteger64Time. Methods defined as returning or accepting a LogicalTimeInterval reference shall return or expect a reference to HLAInteger64TimeInterval (see below).}
\end{compactstdpage}

\clearpage

% --- Standard page 292 / source PDF page 306 ---
\begin{compactstdpage}
    \stdliteraldashitem{An implementation of LogicalTimeFactory called HLAInteger64TimeFactory. All of the methods of this implementation shall return references to HLAInteger64Time, as defined above.}
    \stdliteraldashitem{An implementation of LogicalTimeInterval called HLAInteger64TimeInterval. This implementation shall use an HLAInteger64BE as its underlying representation of a time interval. Methods defined as returning or accepting a LogicalTimeInterval reference shall return or expect a reference to HLAInteger64TimeInterval. The value of epsilon for the HLAInteger64TimeInterval shall be represented by 1. When one of the time types is asked to be incremented by epsilon, the implementation shall increment the value by the most appropriate value to create a value that can be differentiated.}
    \stdliteraldashitem{An implementation of LogicalTimeIntervalFactory called HLAInteger64TimeIntervalFactory. All of the methods of this implementation shall return references to HLAInteger64TimeInterval as defined above.}

    \subsubsection{HLAFloat64Time}

    The implementation of HLAFloat64Time shall provide the following:

    \stdliteraldashitem{An implementation of the interface LogicalTime called HLAFloat64Time. This implementation shall use an HLAFloat64BE as its underlying representation of time. It shall use 0.0 as the value of ``initialTime'' and the largest finite positive value defined for an HLAFloat64BE as its final time. Methods that are defined to return or accept a LogicalTime reference shall actually return or expect a reference to HLAFloat64Time. Methods defined as returning or accepting a LogicalTimeInterval reference shall return or expect a reference to HLAFloat64TimeInterval (see below).}
    \stdliteraldashitem{An implementation of LogicalTimeFactory called HLAFloat64TimeFactory. All of the methods of this implementation shall return references to HLAFloat64Time, as defined above.}
    \stdliteraldashitem{An implementation of LogicalTimeInterval called HLAFloat64TimeInterval. This implementation shall use a HLAFloat64BE as its underlying representation of a time interval. Methods defined as returning or accepting a LogicalTimeInterval reference shall return or expect a reference to HLAFloat64TimeInterval. The value of epsilon for the HLAFloat64TimeInterval is a constant; however, it may result in a larger change when used in an addition or subtraction operation. Since a floating point arithmetic operation may require rounding, the result of adding or subtracting epsilon to a HLAFloat64Time value may round back to the original value. To prevent this conflict with the statements given in 12.3, the implementation of epsilon in an arithmetic operation shall act as the smallest value that guarantees a change in the result. A side effect of this adjustment to epsilon during an arithmetic operation is that the operation may not be reversible (e.g., given T' = T + Ie =\textgreater{} T' -- T \textgreater{}= Ie).}
    \stdliteraldashitem{An implementation of LogicalTimeIntervalFactory called HLAInteger64TimeFactory. All of the methods of this implementation shall return references to HLAInteger64TimeInterval as defined above.}

    \subsection{Connect}

    For all of the APIs, each RTIambassador/FederateAmbassador pair may be part of only one connection to the RTI.

    \subsection{Concurrency and reentrancy}

    The following requirements apply to an RTI:

    \stdliteraldashitem{The RTI shall not make concurrent RTI-initiated service invocations to the same FederateAmbassador instance.}
\end{compactstdpage}

\clearpage

% --- Standard page 293 / source PDF page 307 ---
\begin{compactstdpage}
    \stdliteraldashitem{The RTI shall support concurrent joined federate-initiated service invocations to the same RTIambassador instance. The service invocations shall be processed, in parallel or sequentially.}
    \stdliteraldashitem{When in EVOKED mode, all RTI-initiated service invocations shall be made using the same thread context that was used in the invocation of the Evoke Callback/Evoke Multiple Callbacks service.}
    \stdliteraldashitem{The RTI shall support concurrent joined federate-initiated service invocations to different RTIambassador instances. The calls shall be processed, in parallel or sequentially.}
    \stdliteraldashitem{An RTIambassador instance shall be capable of creating concurrent federation executions.}
    \stdliteraldashitem{The RTI shall support invocation by a joined federate of a subset of joined federate-initiated services from within the context of an RTI-initiated service invocation on the same joined federate. Invocation of joined federate-initiated services not within this subset shall generate the Call Not Allowed From Within Callback exception. The allowed subset of joined federate-initiated services is all joined federate-initiated services except Connect, Disconnect, Join Federation Execution, Resign Federation Execution, Evoke Callback, and Evoke Multiple Callbacks.}

    \subsection{Dynamic link compatibility}

    The APIs defined in this specification were designed to meet the purpose and objectives for dynamic link compatibility of RTIs stated in 1.2 and 1.3 of SISO-STD-004.1-2004 [B5]. The precise details of dynamic link compatibility are language specific and are described in detail in the appropriate language section.

    \subsection{Reflect Attribute Values service methods}

    The Reflect Attribute Values service corresponds to three overloaded versions of reflectAttributeValues() in the Java and C++ APIs. Which version the RTI calls is dependent on if the updateAttributeValues() was called with a timestamp and if the message was received TSO. The optional arguments are encapsulated in a SupplementalReflectInfo object that is added as an argument to the Reflect Attribute Values service method. The values of the convey region designator switch and the Convey Producing Federate Switch control the contents of the SupplementalReflectInfo object.

    The first version passes no optional arguments but does pass a SupplementalReflectInfo object with possibly valid values. This version is called when there is not a timestamp associated with the updateAttributeValues() call and the message was received RO. The second version passes the timestamp, receive message order type optional arguments, and a SupplementalReflectInfo object with possibly valid values. This version is called when there is a timestamp associated with the updateAttributeValues() call and the message was received RO. The third version passes the timestamp, receive message order type, message retraction designator optional arguments, and a SupplementalReflectInfo object with possibly valid values. This version is called when there is a timestamp associated with the updateAttributeValues() call and the message was received TSO.

    \subsection{Receive Interaction service methods}

    The Receive Interaction service corresponds to three overloaded versions of receiveInteraction() in the Java and C++ APIs. Which version the RTI calls is dependent on if the sendInteraction() was called with a timestamp and if the message was received TSO. The optional arguments are encapsulated in a SupplementalReceiveInfo object that is added as an argument to the Receive Interaction service method. The values of the convey region designator switch and the Convey Producing Federate Switch control the contents of the SupplementalReceiveInfo object.

    The first version passes no optional arguments but does pass a SupplementalReceiveInfo object with possibly valid values. This version is called when there is not a timestamp associated with the sendInteraction() call and the message was received RO. The second version passes the timestamp,

\end{compactstdpage}

\clearpage

% --- Standard page 294 / source PDF page 308 ---
\begin{compactstdpage}
    receive message order type optional arguments, and a SupplementalReceiveInfo object with possibly valid values. This version is called when there is a timestamp associated with the sendInteraction() call and the message was received RO. The third version passes the timestamp, receive message order type, message retraction designator optional arguments, and a SupplementalReceiveInfo object with possibly valid values. This version is called when there is a timestamp associated with the sendInteraction() call and the message was received TSO.

    \subsection{Remove Object Instance service methods}

    The Remove Object Instance service corresponds to three overloaded versions of the removeObjectInstance() in the Java and C++ APIs. Which version the RTI calls is dependent on if the removeObjectInstance() was called with a timestamp and if the message was received TSO. The optional arguments are in a SupplementalRemoveInfo object that is added as an argument to the Remove Object Instance service method. The values of the Convey Producing Federate Switch control the contents of the SupplementalReceiveInfo object.

    The first version returns no optional arguments but does pass a SupplementalRemoveInfo object with possibly valid values. This version is called when there is not a timestamp associated with the removeObjectInstance() call and the message was received RO. The second version returns the timestamp, receive message order type optional arguments, and a SupplementalRemoveInfo object with possibly valid values. This version is called when there is a timestamp associated with the removeObjectInstance() call and the message was received RO. The third version returns the timestamp, receive message order type, message retraction designator optional arguments, and a SupplementalRemoveInfo object with possibly valid values. This version is called when there is a timestamp associated with the removeObjectInstance() call and the message was received TSO.

    \subsection{Java}

    The Java API is composed of interfaces and classes. Where an interface is specified, the RTI implementers shall furnish an implementation of the interface. For instance, RTIambassador is specified as an interface. Therefore, an RTI implementer shall furnish an implementation of the interface, for instance:

        {\ttfamily\fontsize{8.7}{10}\selectfont package com.myrticompany.amazingrti;\\
    public final class MyRTIimpl implements hla.rti1516e.RTIambassador \{\ldots\}}

    Each such interface as specified remains part of the API. Thus, federate code can use the specified interface type. The FederateAmbassador, however, is an exception to this rule and shall be implemented by a federate developer.

    Classes defined in the Java API, such as RTIexception and ResignAction, shall be included in each implementation of the RTI.

    The entire Java API is included in the package hla.rti1516e and subpackages. Because the specified interfaces and classes constitute the API, the entire package hla.rti1516e shall be furnished without additions or deletions or change of package name as part of each implementation. Implementers shall place their implementation classes in another package.

    \subsubsection{Designators}

    The Java interfaces that represent handles shall implement a hashCode() method that returns an int that is suitable for use as a hash value for use in storing handles in a hash table. The hash value may not be unique across two separate handle values but will be equal given two separate instances of the same handle

\end{compactstdpage}

\clearpage

% --- Standard page 295 / source PDF page 309 ---
\begin{compactstdpage}
    value. The following are the properties of the hash value where H1 and H2 are separate instances of hash values.

    \stdliteraldashitem{H1 == H2 implies H1.hashCode() == H2.hashCode()}
    \stdliteraldashitem{H1 != H2 does not imply H1.hashCode() != H2.hashCode()}

    \subsubsection{Dynamic link compatibility}

    Java dynamic link compatibility for this specification is provided through use of the service registry, as described in the javadoc for ServiceRegistry, available at \url{http://java.sun.com/j2se/1.5.0/docs/api/javax/imageio/spi/ServiceRegistry.html}.

    The RTI shall provide the HLA 1516-2009.1 Java API classes in the package hla.rti1516e with corresponding subpackages. Each RTI vendor shall provide their RTI implementation classes in an implementation-specific package.

    This standard provides implementations for the classes RtiFactoryFactory, NullFederateAmbassador, and LogicalTimeFactoryFactory.

    The RTI shall provide implementations for the Encoding Helper classes in the hla.rti1516e.encoding package. The RTI shall also provide a class that implements the RtiFactory interface defined in the hla.rti1516e package. This class shall reside in an implementation-specific package outside of the hla.rti1516e hierarchy. The RtiFactory class shall be listed in a service registry entry, i.e., a file named hla.rti1516e.RtiFactory in the META-INF/services directory of the Java Archive (JAR) shall contain a line with the class name of the implementing class, e.g., com.mycompany.MyRtiFactoryImpl.

    {\sloppy The RTI shall also provide implementations of the standardized time classes HLAinteger64Time,\allowbreak{} HLAinteger64Interval,\allowbreak{} HLAfloat64Time,\allowbreak{} and HLAfloat64Interval,\allowbreak{} along with their factory classes HLAinteger64TimeFactory and HLAfloat64TimeFactory. The time factories shall be listed in a service registry entry, i.e., a file named hla.rti1516e.LogicalTimeFactory in the META-INF/services directory of the JAR shall contain a line with the class name of each implementing class, e.g., com.mycompany.HLAinteger64TimeFactoryImpl\allowbreak{} and com.mycompany.HLAfloat64Time-FactoryImpl.\par}

    All RTIs shall support the interpretation of the FDD designator as a URL.

    \paragraph{Factory mechanism}

    An RTI implementation shall furnish an implementation of the hla.rti1516e.RtiFactory interface.

    \subparagraph{Registration}

    An RTI implementation shall register as a service by creating a file in the services subdirectory of the META-INF directory in the JAR file containing the RTI implementation. The file shall be named hla.rti1516e.RtiFactory and contain a line specifying the name of the class that implements the RtiFactory interface, e.g.,

            {\ttfamily\fontsize{8.7}{10}\selectfont org.someorganization.anrti.RtiFactoryImpl}

    \subparagraph{RtiFactoryFactory}

    To acquire a reference to an RtiFactory implementation, the federate uses the RtiFactoryFactory class provided with the standard.

\end{compactstdpage}

\clearpage

% --- Standard page 296 / source PDF page 310 ---
\begin{compactstdpage}
{\ttfamily\fontsize{8.7}{10}\selectfont
RtiFactory rtiFactory = RtiFactoryFactory.getRtiFactory();\\
RTIambassador \_rtiAmbassador = rtiFactory.createRtiAmbassador();}

    The federate can acquire a named RtiFactory,

        {\ttfamily\fontsize{8.7}{10}\selectfont
    RtiFactory rtiFactory = RtiFactoryFactory.getRtiFactory("SomeRTI");\\
    RTIambassador \_rtiAmbassador = rtiFactory.createRtiAmbassador();}

    or get a list of available RtiFactories.

        {\ttfamily\fontsize{8.7}{10}\selectfont Set\textless{}RtiFactory\textgreater{} rtiFactories = RtiFactoryFactory.getRtiFactories();}

    \subparagraph{RtiFactory}

    The RtiFactory class or the Java API for this specification follows the RtiFactory class description detailed in 6.6.1 of SISO-STD-004.1-2004 [B5].

    In addition, this specification also supplies a getEncoderFactory() method that returns the EncoderFactory that can be used by the encoding services described in 12.11.3.2.

    \subsubsection{Additional functionality}

    Implementations of the Java API shall add the following utility classes to the hla.rti1516e package.

    \paragraph{NullFederateAmbassador class}

    The NullFederateAmbassador class provides empty implementations of all methods in FederateAmbassador.

    \paragraph{encoding package}

    The encoding package provides a set of classes providing helper functions for encoding datatypes according to IEEE Std 1516.2-2010. These classes are in a package named encoding within the hla.rti1516e package (e.g., hla.rti1516e.encoding)

    All encoding helper classes shall support the encode and decode methods for converting language-specific datatypes to and from standard HLA encoded byte-arrays. All noncomplex encoder classes, as listed in Table 19, shall support the getValue and setValue methods for retrieving and updating the value.

    All complex encoder classes shall support elements that make up the complex class. It shall be possible to use both complex and noncomplex encoder class instances as elements of a complex class instance.

    A method named add shall be supported for adding elements to an instance of such a complex class. See Table 20. Another method named get(index) for retrieving data elements shall also be supported.

    The HLAvariantRecord shall support the method setVariant to associate different encoders/decoders with different discriminants and the methods setDiscriminant and getDiscriminant to indicate which variant of the record that shall be used for composing, encoding, and decoding data elements.

    The HLAVariableArray encoder shall provide a size method that provides the number of elements.

    The RTI shall supply an implementation of the interface hla.rti1516e.encoding.EncoderFactory, a factory that shall provide methods for creating encoding helper instances. The RtiFactory shall supply a method getEncoderFactory that returns an EncoderFactory instance.

\end{compactstdpage}

\clearpage

% --- Standard page 297 / source PDF page 311 ---
\begin{compactstdpage}
    \cxiitablenineteenandtwenty

\end{compactstdpage}

\clearpage

% --- Standard page 298 / source PDF page 312 ---
\begin{compactstdpage}
    \subsubsection{Services}

    In general, each service in Clause 4 through Clause 10 corresponds to one method in each API. In some cases, because of optional arguments to a service or programming language constraints, a single service may be mapped to multiple methods in an API.

    For the Java API, the federate-initiated services invoked on the RTI are methods of the RTIambassador interface, and RTI-initiated service callbacks on the federate are methods of the FederateAmbassador interface. The RTI and federate ambassador services in the Java API that do not follow a one-to-one mapping with the services in Clause 4 through Clause 10 of the specification are described in 12.11.4.1 through 12.11.4.34.

    \paragraph{Create Federation Execution service}

    The set of FOM module designators argument is implemented as an array of URLs where each item in the array is a URL for a file containing one FOM module.

    In addition to the supplied arguments listed with the description of this service, the Java API requires the federate to supply a String logicalTimeImplementationName that contains the name of the logical time implementation. In addition to the specified exceptions listed with the description of this service, the Java API may also throw a CouldNotCreateLogicalTimeFactory. This exception indicates that the RTI was not able to create an instance of the logical time factory class that was named in the String logicalTimeImplementationName argument.

    \paragraph{Join Federation Execution service}

    The set of FOM module designators argument is implemented as an array of URLs where each item in the array is a URL for a file containing one FOM module.

    In addition to the specified exceptions listed with the description of this service, the Java API may also throw a CouldNotCreateLogicalTimeFactory. This exception indicates that the RTI was not able to create an instance of the logical time factory class that was named in the String logicalTimeImplementationName argument to the Create Federation Execution service.

    \paragraph{Register Federation Synchronization Point service}

    The optional argument representing the set of joined federate designators is implemented in the Java API as two registerFederationSynchronizationPoint() methods: one that takes a FederateHandleSet as an argument, and one that does not.

    \paragraph{Confirm Synchronization Point Registration \rtiinitiated{} service}

    The registration-success indicator is implemented in the Java API as two distinct methods: synchronizationPointRegistrationSucceeded()                and synchronizationPointRegistrationFailed(). A Synchronization Point Failure Reason (discussed in 12.11.5.10), in the form of a SynchronizationPointFailureReason, is passed as an argument to synchronizationPointRegistrationFailed().

    \paragraph{Synchronization Point Achieved service}

    The optional synchronization-success indicator parameter is implemented as two methods, one of which takes a boolean as an argument.

\end{compactstdpage}

\clearpage

% --- Standard page 299 / source PDF page 313 ---
\begin{compactstdpage}
    \paragraph{Request Federation Save service}

    The optional timestamp argument is implemented as two methods, one of which takes a LogicalTime as an argument.

    \paragraph{Initiate Federate Save \rtiinitiated{} service}

    The optional timestamp argument is implemented as two methods, one of which takes a LogicalTime as an argument.

    \paragraph{Federate Save Complete service}

    This service corresponds to two methods in the Java API: federateSaveComplete() and federateSaveNotComplete().

    \paragraph{Federation Saved \rtiinitiated{} service}

    This service corresponds to two methods in the Java API: federationSaved() and federationNotSaved(). A Save Failure Reason (discussed in 12.11.5.11), in the form of a SaveFailureReason, is passed as an argument to federationNotSaved().

    \paragraph{Confirm Federation Restoration Request \rtiinitiated{} service}

    This service is mapped to the following two methods in the Java API: requestFederationRestoreSucceeded() and requestFederationRestoreFailed().

    \paragraph{Federate Restore Complete service}

    This service is mapped to the following two methods in the Java API: federateRestoreComplete() and federateRestoreNotComplete().

    \paragraph{Federation Restored \rtiinitiated{} service}

    This service corresponds to two methods in the Java API: federationRestored() and federationNotRestored(). A Restore Failure Reason (discussed in 12.11.5.12), in the form of a RestoreFailureReason, is passed as an argument to federationNotRestored().

    \paragraph{Connect service}

    In addition to the supplied arguments listed with the description of this service, the Java API requires the federate to supply an instance of the interface FederateAmbassador. The optional localSettingsDesignator parameter is required and can contain RTI implementation-specific configuration information. An empty string for the localSettingsDesignator parameter shall be used to indicate the default local settings.

    \paragraph{Unpublish Object Class Attributes service}

    The optional argument representing the set of attribute designators is implemented in the Java API as two methods: unpublishObjectClass() and unpublishObjectClassAttributes(). The latter of the two takes an AttributeHandleSet as an argument.

\end{compactstdpage}

\clearpage

% --- Standard page 300 / source PDF page 314 ---
\begin{compactstdpage}
    \paragraph{Subscribe Object Class Attributes service}

    The optional passive subscription indicator and optional update rate designator are implemented in the Java API as two overloaded methods: subscribeObjectClassAttributes() and subscribeObjectClassAttributesPassively(). One version of each method has only the required arguments, and the second version adds the optional String updateRateDesignator argument.

    \paragraph{Unsubscribe Object Class Attributes service}

    The optional argument representing the set of attribute designators is implemented in the Java API as two methods: unsubscribeObjectClass() and unsubscribeObjectClassAttributes(). The latter of the two takes an AttributeHandleSet as an argument.

    \paragraph{Subscribe Interaction Class service}

    The optional passive subscription indicator is implemented in the Java API as two methods: subscribeInteractionClass() and subscribeInteractionClassPassively().

    \paragraph{Reserve Object Instance Name service}

    The optional single name parameter or list of names parameter is implemented in the Java API as two methods: reserveObjectInstanceName() and reserveMultipleObjectInstanceName().

    \paragraph{Object Instance Name Reserved \rtiinitiated{} service}

    This service corresponds to two methods in the Java API: objectInstanceNameReservationSucceeded() and objectInstanceNameReservationFailed().

    \paragraph{Multiple Object Instance Name Reserved \rtiinitiated{} service}

    This service corresponds to two methods in the Java API: multipleObjectInstanceNameReservationSucceeded() and multipleObjectInstanceNameReservationFailed().

    \paragraph{Register Object Instance service}

    The optional argument representing the object instance name is implemented in the Java API as two registerObjectInstance() methods: one that takes a String as an argument for the object instance name, and one that does not.

    \paragraph{Discover Object Instance service}

    The optional argument representing the producing joined federate designator is implemented in the Java API as two discoverObjectInstance() methods: one that takes a FederateHandle as an argument for the producing joined federate designator, and one that does not.

    \paragraph{Update Attribute Values service}

    The optional argument representing the timestamp is implemented in the Java API as two updateAttributeValues() methods: one that takes a LogicalTime as an argument for the timestamp, and one that does not. The version that takes a LogicalTime also returns a MessageRetractionReturn for the optional message retraction designator.

\end{compactstdpage}

\clearpage

% --- Standard page 301 / source PDF page 315 ---
\begin{compactstdpage}
    \paragraph{Reflect Attribute Values \rtiinitiated{} service}

    This service corresponds to three overloaded versions of the reflectAttributeValues() method in the Java API. Which version the RTI calls is dependent on if the updateAttributeValues() was called with a timestamp and if the message was received TSO. See 12.8 for details how to determine which version the RTI will call.

    \paragraph{Send Interaction service}

    The optional argument representing the timestamp is implemented in the Java API as two sendInteraction() methods: one that takes a LogicalTime as an argument for the timestamp, and one that does not. The version that takes a LogicalTime also returns a MessageRetractionReturn for the optional message retraction designator.

    \paragraph{Receive Interaction \rtiinitiated{} service}

    This service corresponds to three overloaded versions of the receiveInteraction() method in the Java API. Which version the RTI calls is dependent on if the sendInteraction() was called with a timestamp and if the message was received TSO. See 12.9 for details on how to determine which version the RTI will call.

    \paragraph{Delete Object Instance service}

    The optional argument representing the timestamp is implemented in the Java API as two deleteObjectInstance() methods: one that takes a LogicalTime as an argument for the timestamp, and one that does not. The version that takes a LogicalTime also returns a MessageRetractionReturn for the optional message retraction designator.

    \paragraph{Remove Object Instance \rtiinitiated{} service}

    This service corresponds to three overloaded versions of the removeObjectInstance() method in the Java API. See 12.10 for details on how to determine which version the RTI will call.

    \paragraph{Request Attribute Value Update service}

    This service corresponds to two overloaded methods named requestAttributeValueUpdate() in the Java API. One takes an ObjectInstanceHandle to represent the object instance designator. The other takes an ObjectClassHandle to represent the object class designator.

    \paragraph{Inform Attribute Ownership \rtiinitiated{} service}

    In the Java API, this service corresponds to three distinct methods: attributeIsNotOwned(), attributeIsOwnedByRTI(), and informAttributeOwnership(). All three methods take an ObjectInstanceHandle to represent the object instance designator and an AttributeHandle to represent the attribute designator as arguments. In addition, informAttributeOwnership() takes a FederateHandle as an argument to represent the federate designator.

    \paragraph{Register Object Instance With Regions service}

    This service corresponds to two overloaded versions of the registerObjectInstanceWithRegions() method in the Java API. The second version takes a String as an argument to represent the optional object instance name.

\end{compactstdpage}

\clearpage

% --- Standard page 302 / source PDF page 316 ---
\begin{compactstdpage}
    \paragraph{Subscribe Object Class Attributes With Regions service}

    The optional passive subscription indicator and optional update rate designator are implemented in the Java API as two overloaded methods: subscribeObjectClassAttributesWithRegions() and subscribeObjectClassAttributesPassivelyWithRegions(). One version of each method has only the required arguments, and the second version adds the optional String updateRateDesignator argument.

    \paragraph{Subscribe Interaction Class With Regions service}

    The optional passive subscription indicator is implemented in the Java API as two methods: subscribeInteractionClassWithRegions()               and subscribeInteractionClassPassivelyWithRegions().

    \paragraph{Send Interaction With Regions service}

    The optional argument representing the timestamp is implemented in the Java API as two sendInteractionWithRegions() methods: one that takes a LogicalTime as an argument for the timestamp, and one that does not. The version that takes a LogicalTime also returns a MessageRetractionReturn for the optional message retraction designator.

    \subsubsection{Abstract datatypes}

    In 1.4.2, general conventions used to describe the abstract datatypes used in Clause 4 through Clause 10 are presented. In 12.11.5.1 through 12.11.5.26, the abstract datatypes are further detailed by mapping an individual datatype to its implementation in the Java API. Note that the datatypes from Clause 11 are explained in that clause and are not repeated here. Additionally, designators are a subject in their own right and are covered in 12.2.

    Abstract datatypes are mapped to idiomatic Java types (if possible, to a type in the Java Platform APIs). Java 5 is assumed. This is especially evident in the use of the Java Development Kit Collections framework. All interfaces extend, and classes implement, java.io.Serializable. The abstract datatypes make use of generics (typed collections) to achieve type safety in collections.

    \paragraph{Names, labels, and strings}

    The names, labels, and strings described in Clause 4 through Clause 10 of this specification are mapped to java.lang.String.

    \paragraph{Boolean values}

    The Java primitive type boolean is used to represent all boolean values.

    \paragraph{Exceptions}

    All exceptions in the Java API inherit from a class named RTIexception, which extends java.lang.Exception. This hierarchy allows all exceptions in the Java API to be caught as a single base class.

    \paragraph{Handles}

    Most designators in the Java API are handles. Each handle is a distinct Java interface that shall be implemented by the RTI implementer. Each handle interface supports equals() and hashCode() so that handles can be used as keys in HashTables. Handles of different types are incommensurable.

\end{compactstdpage}

\clearpage

% --- Standard page 303 / source PDF page 317 ---
\begin{compactstdpage}
    Handle instances typically are returned from RTI services, e.g., getAttributeHandle(). These services return a reference to the handle interface and, thus, hide the implementation class. The federate programmer typically uses only handle references that have been returned by the RTI.

    Each handle interface supports encode(), which shall place a compact, opaque representation suitable for network transmission into the provided byte array at the specified offset. Each handle interface has a corresponding handle factory interface, which shall also be implemented by the RTI implementer. The factory interface defines decode(), which returns a new handle instance from a representation in the provided byte array at the indicated offset. This mechanism is intended to allow handles to be saved and restored later as well as to allow handles to be passed among federates as attribute values or parameter values.

    Handle classes also provide a toString() method that returns a printable form of a handle. The string produced, however, is not guaranteed to be the same as the name to which the handle corresponds (assuming there is a name that corresponds to that kind of handle). If the user needs the name form of a handle, employ the appropriate support service from Clause 10.

    \paragraph{Resignation directive}

    The Resign Federation Execution service allows the specification of one of six directives, which are represented by an enumeration, ResignAction. The six different directives of the Resign Federation Execution service are defined as the six enumerated values that instances of the ResignAction enumeration may take. The ResignAction enumeration has the following values:

    \stdliteraldashitem{UNCONDITIONALLY\_DIVEST\_ATTRIBUTES}
    \stdliteraldashitem{DELETE\_OBJECTS}
    \stdliteraldashitem{CANCEL\_PENDING\_OWNERSHIP\_ACQUISITIONS}
    \stdliteraldashitem{DELETE\_OBJECTS\_THEN\_DIVEST}
    \stdliteraldashitem{CANCEL\_THEN\_DELETE\_THEN\_DIVEST}
    \stdliteraldashitem{NO\_ACTION}

    \paragraph{Attribute and parameter values and user-supplied tag}

    Several services in the HLA interface allow the federate programmer to supply additional information to be transported by the RTI in an opaque manner, such as the Update Attribute Values service. Attribute and parameter values are represented as byte arrays, byte[]. byte[] is also used to supply or deliver user-supplied tags.

    \paragraph{Sets of designators}

    Sets of designators, e.g., federate designators in the Register Federation Synchronization Point service or attribute designator sets in the Publish Object Class Attributes service, are represented by interfaces that extend java.util.Set for the respective handle type. This interface is part of the Java 2 Collections framework. The RTI implementer's implementation of each such set shall adhere to the following:

    \stdliteraldashitem{Members shall be of the appropriate type, e.g., AttributeHandle for AttributeHandleSet. Attempts to add an instance of an inappropriate type shall elicit IllegalArgumentException.}
    \stdliteraldashitem{Implementations of the sets shall implement all defined operations (none are optional).}
    \stdliteraldashitem{Methods that might take an incommensurable set as an actual parameter shall throw IllegalArgumentException.}
\end{compactstdpage}

\clearpage

% --- Standard page 304 / source PDF page 318 ---
\begin{compactstdpage}
    Each such handle set has a corresponding factory interface. Thus, the implementation class shall be hidden from the federate programmer. A set of handles is provided for federate handles, attribute handles, dimension handles, and region handles.

    \paragraph{Success indicators}

    Several of the requests a federate may make of an RTI shall be acknowledged by the RTI to the federate. This acknowledgment can indicate that the request succeeded or failed. In the Java API, each request acknowledgment corresponds to a positive and a negative callback method on the FederateAmbassador, e.g., federationSaved() and federationNotSaved(), requestFederationRestoreSucceeded() and requestFederationRestoreFailed(), and federationRestored() and federationNotRestored().

    \paragraph{Callback Model}

    When the connect() method is invoked by a federate, it shall supply an instance of the CallbackModel enumeration as an argument. The CallbackModel enumeration provides the following values:

    \stdliteraldashitem{HLA\_IMMEDIATE}
    \stdliteraldashitem{HLA\_EVOKED}

    \paragraph{Synchronization Point Failure Reason}

    When the synchronizationPointRegistrationFailed() method is invoked by the RTI on a federate, the RTI provides an instance the SynchronizationPointFailureReason enumeration as an argument. The SynchronizationPointFailureReason enumeration has the following values:

    \stdliteraldashitem{SYNCHRONIZATION\_POINT\_LABEL\_NOT\_UNIQUE}
    \stdliteraldashitem{SYNCHRONIZATION\_SET\_MEMBER\_NOT\_JOINED}

    \paragraph{Save Failure Reason}

    When the federationNotSaved() method is invoked by the RTI on a federate, the RTI provides an instance of the SaveFailureReason enumeration as an argument. The SaveFailureReason enumeration has the following values:

    \stdliteraldashitem{RTI\_UNABLE\_TO\_SAVE}
    \stdliteraldashitem{FEDERATE\_REPORTED\_FAILURE\_DURING\_SAVE}
    \stdliteraldashitem{FEDERATE\_RESIGNED\_DURING\_SAVE}
    \stdliteraldashitem{RTI\_DETECTED\_FAILURE\_DURING\_SAVE}
    \stdliteraldashitem{SAVE\_TIME\_CANNOT\_BE\_HONORED}
    \stdliteraldashitem{SAVE\_ABORTED}

    \paragraph{Restore Failure Reason}

    When the federationNotRestored() method is invoked by the RTI on a federate, the RTI provides an instance of the RestoreFailureReason enumeration as an argument. The RestoreFailureReason enumeration has the following values:

    \stdliteraldashitem{RTI\_UNABLE\_TO\_RESTORE}
    \stdliteraldashitem{FEDERATE\_REPORTED\_FAILURE\_DURING\_RESTORE}
    \stdliteraldashitem{FEDERATE\_RESIGNED\_DURING\_RESTORE}
    \stdliteraldashitem{RTI\_DETECTED\_FAILURE\_DURING\_RESTORE}
\end{compactstdpage}

\clearpage

% --- Standard page 305 / source PDF page 319 ---
\begin{compactstdpage}
    \stdliteraldashitem{RESTORE\_ABORTED}

    \paragraph{List of joined federates save status}

    The list of joined federates and the save status of each is implemented as an array of FederateHandleSaveStatusPair objects. Each FederateHandleSaveStatusPair object contains a FederateHandle and a SaveStatus representing the save status for that joined federate. The SaveStatus enumeration has the following values:

    \stdliteraldashitem{NO\_SAVE\_IN\_PROGRESS}
    \stdliteraldashitem{FEDERATE\_INSTRUCTED\_TO\_SAVE}
    \stdliteraldashitem{FEDERATE\_SAVING}
    \stdliteraldashitem{FEDERATE\_WAITING\_FOR\_FEDERATION\_TO\_SAVE}

    \paragraph{List of joined federates restore status}

    The list of joined federates and the restore status of each is implemented as an array of FederateRestoreStatus objects. Each FederateRestoreStatus object contains a FederateHandle that contains the pre-restore federate handle, a FederateHandle that contains the post-restore federate handle, and a RestoreStatus representing the restore status for that joined federate. The RestoreStatus enumeration has the following values:

    \stdliteraldashitem{NO\_RESTORE\_IN\_PROGRESS}
    \stdliteraldashitem{FEDERATE\_RESTORE\_REQUEST\_PENDING}
    \stdliteraldashitem{FEDERATE\_WAITING\_FOR\_RESTORE\_TO\_BEGIN}
    \stdliteraldashitem{FEDERATE\_PREPARED\_TO\_RESTORE}
    \stdliteraldashitem{FEDERATE\_RESTORING}
    \stdliteraldashitem{FEDERATE\_WAITING\_FOR\_FEDERATION\_TO\_RESTORE}

    \paragraph{Passive subscription indicator}

    The Subscribe Object Class Attributes, Subscribe Interaction Class, Subscribe Object Class Attributes With Regions, and Subscribe Interaction Class With Regions services each may optionally provide a boolean value to indicate whether the subscription is active. In the Java API, this corresponds to a pair of methods on RTIambassador,                 e.g.,           subscribeObjectClassAttributes()                            and subscribeObjectClassAttributesPassively().

    \paragraph{Constrained set of attribute designator and value pairs}

    The Update Attribute Values and Reflect Attribute Values \rtiinitiated{} services require as arguments a set of attribute designator and value pairs. Such sets are instances of AttributeHandleValueMap. This interface extends java.util.Map, part of the Java 2 Collections framework. The RTI implementer's implementation of AttributeHandleValueMap shall adhere to the following:

    \stdliteraldashitem{The keys are instances of AttributeHandle; IllegalArgumentException shall be thrown for violations.}
    \stdliteraldashitem{The values are instances of byte[]; IllegalArgumentException shall be thrown for violations.}
    \stdliteraldashitem{All java.util.Map operations shall be implemented; none are optional.}
    \stdliteraldashitem{The implementation shall not accept null mappings, i.e., a mapping to a null value.}
    The interface AttributeHandleValueMapFactory shall be implemented, and the implementation class for AttributeHandleValueMap shall be hidden from the federate programmer. The federate

\end{compactstdpage}

\clearpage

% --- Standard page 306 / source PDF page 320 ---
\begin{compactstdpage}
    programmer shall obtain a reference to a factory instance by the method getAttributeHandleValueMapFactory() on RTIambassador.

    \paragraph{Constrained set of interaction parameter designator and value pairs}

    The Send Interaction, Receive Interaction \rtiinitiated{}, and Send Interaction With Regions services require as arguments a constrained set of parameter designator and value pairs. Such sets are instances of ParameterHandleValueMap. This interface extends java.util.Map, part of the Java 2 Collections framework. The RTI implementer's implementation of ParameterHandleValueMap shall adhere to the following:

    \stdliteraldashitem{The keys are instances of ParameterHandle; IllegalArgumentException shall be thrown for violations.}
    \stdliteraldashitem{The values are instances of byte[]; IllegalArgumentException shall be thrown for violations.}
    \stdliteraldashitem{All java.util.Map operations shall be implemented; none are optional.}
    \stdliteraldashitem{The implementation shall not accept null mappings, i.e., a mapping to a null value.}
    The interface ParameterHandleValueMapFactory shall be implemented, and the implementation class for ParameterHandleValueMap shall be hidden from the federate programmer. The RTI shall provide the federate programmer a reference to a factory instance by the method getParameterHandleValueMapFactory() on RTIambassador.

    \paragraph{Message order type}

    Message order types are represented by an enumerated type. The following enumerations are provided: RECEIVE and TIMESTAMP. These enumerations can be passed to the RTI or compared to an instance of an OrderType object returned by the RTI.

    The OrderType enumeration shares some aspects of the interface to handles (see 12.11.5.4); specifically, the means to encode an order type, reconstitute an encoded order type, and produce a printable form of an order type are provided. Encoding order types and reconstitution of encoded order types work in the same way as for handles, with the exception that the decode() method is provided as a static member of the OrderType enumerated type rather than as a member of a corresponding factory. Producing a printable form works the same as for handles. The toString() method is used, and the string produced is not guaranteed to be the same as the name to which the order type corresponds (i.e., the name of the order type as found in the FDD). If the user needs the name form of an order type, employ the Get Order Name service.

    \paragraph{Transportation type}

    The Java API uses a handle, TransportationTypeHandle, to represent all of the arguments of type TransportationType in the API methods. To acquire a handle, use the Get Transportation Type Handle service. All RTIs shall support the transportation types HLA\_RELIABLE and HLA\_BEST\_EFFORT.

    It is possible that some implementations of the RTI will allow additional transportation types. In this case, those RTI implementations shall accept the new transportation type name in the FDD and the implementation shall

    \stdliteraldashitem{Return a valid Transportation Type Handle for the Get Transportation Type Handle service.}
    \stdliteraldashitem{Return a valid Transportation Type Name for the Get Transportation Type Name service.}
\end{compactstdpage}

\clearpage

% --- Standard page 307 / source PDF page 321 ---
\begin{compactstdpage}
    \paragraph{Ownership designator}

    The Inform Attribute Ownership \rtiinitiated{} service supplies an ownership designator as one of its arguments. In the Java API, this corresponds to three distinct methods: attributeIsNotOwned(), attributeIsOwnedByRTI(), and informAttributeOwnership(). All three methods take an ObjectInstanceHandle and an AttributeHandle as arguments. In addition, informAttributeOwnership() takes a FederateHandle as an argument. Together, these three methods combine to represent the ownership designator of the Inform Attribute Ownership \rtiinitiated{} service.

    \paragraph{Definition indicator}

    The Query GALT and Query LITS services are required to return an indicator to define whether their return value is valid. Both services return a value of type TimeQueryReturn. This structure contains a boolean member that indicates whether the LogicalTime is valid. If the member timeIsValid is false, the member time should not be used. TimeQueryReturn also supports equals() and hashCode() so that TimeQueryReturn objects can be used as keys in HashTables.

    \paragraph{Optional message retraction designator}

    The Update Attribute Values, Send Interaction, Send Interaction With Regions, and Delete Object Instance services may return a message retraction designator under certain circumstances. In the Java API, this indicator is implemented much like the definition indicator described above. The methods corresponding to these services always return a MessageRetractionReturn. The member handle of this structure should be used only if the member retractionHandleIsValid is true.

    \paragraph{Collection of attribute designator set and region designator set pairs}

    Several services require the federate to provide the RTI with a collection of attribute designator set and region designator set pairs. In the Java API, this set corresponds to the interface AttributeSetRegionSetPairList. This interface extends java.util.List, part of the Java 2 Collections framework. The RTI implementer's implementation of AttributeSetRegionSetPairList shall adhere to the following:

    \stdliteraldashitem{The elements are instances of the class AttributeRegionAssociation; IllegalArgumentException shall be thrown for violations.}
    \stdliteraldashitem{All java.util.List operations shall be implemented; none are optional.}
    \stdliteraldashitem{The implementation shall not accept null elements; i.e., an element whose value is null.}
    \stdliteraldashitem{No operation shall accept an incommensurate type; IllegalArgumentException shall be thrown for violations.}

    \paragraph{Logical time, timestamps, and lookahead}

    In the Java API, the abstract interfaces for the logical time classes are defined as Java interfaces. The LogicalTime and LogicalTimeInterval interfaces extend the Comparable interface. LogicalTime and LogicalTimeInterval interfaces shall have a toString() method, which yields the representation suitable for printing. The variable array used with the encoding and decoding methods shall be a byte array.

    In addition to the methods mentioned in 12.3, the Java LogicalTime interface provides the following methods:

    \stdliteraldashitem{hashCode() that returns an integer value that shall be equal for any two LogicalTime instances for which equals() returns true.}
    \stdliteraldashitem{encodedLength() that shall return the size of the buffer required to encode the object}
\end{compactstdpage}

\clearpage

% --- Standard page 308 / source PDF page 322 ---
\begin{compactstdpage}
    \stdliteraldashitem{encode() that takes a byte array and an offset and shall encode the object into the byte array at the offset position}
    In addition to the methods mentioned in 12.3, the Java LogicalTimeInterval interface provides the following methods:

    \stdliteraldashitem{hashCode() that returns an integer value that shall be equal for any two LogicalTime instances for which equals() returns true.}
    \stdliteraldashitem{encodedLength() that shall return the size of the buffer required to encode the object}
    \stdliteraldashitem{encode() that takes a byte array and an offset and shall encode the object into the byte array at the offset position}
    In addition to the methods mentioned in 12.3, the Java LogicalTimeFactory interface provides the following methods:

    \stdliteraldashitem{decodeTime() that takes a byte array with an offset and returns a LogicalTime instance with the value that is decoded from that offset in the byte array. This method can throw a CouldNotDecode exception.}
    \stdliteraldashitem{decodeInterval() that takes a byte array with an offset and returns a LogicalTimeInterval instance with the value that is decoded from that offset in the byte array. This method can throw a CouldNotDecode exception.}
    \stdliteraldashitem{getName() that returns the String name of the logical time implementation that was passed to the Create Federation Execution service.}

    \subparagraph{Registration}

    A time implementation shall register as a service by creating a file in the services subdirectory of the META-INF directory in the JAR file containing the time implementation. The file shall be named hla.rti1516e.LogicalTimeFactory and contain a line specifying the name of the class that implements the LogicalTimeFactory interface, e.g.,

        {\ttfamily\fontsize{8.7}{10}\selectfont org.someorganization.time.HLAinteger64TimeFactoryImpl}

    \subparagraph{Standardized time types}

    An RTI implementation shall provide implementations of the standardized time types as classes implementing the interfaces HLAinteger64Time, HLAinteger64Interval, HLAinteger64TimeFactory, HLAfloat64Time, HLAfloat64Interval, and HLAfloat64TimeFactory. The getName() methods of the HLAinteger64TimeFactory and HLAfloat64TimeFactory implementations shall return the symbolic names ``HLAinteger64Time'' and ``HLAfloat64Time,'' respectively. The HLAinteger64TimeFactory and HLAfloat64TimeFactory implementations shall be registered according to 12.11.5.24.1.

    \subparagraph{LogicalTimeFactoryFactory}

    To acquire a LogicalTimeFactory implementation, the federate may use the LogicalTimeFactoryFactory class, provided with this standard, and supply the symbolic name,

        {\ttfamily\fontsize{8.3}{9.6}\selectfont
    HLAinteger64TimeFactory logicalTimeFactory = (HLAinteger64TimeFactory)\\
    LogicalTimeFactoryFactory.getLogicalTimeFactory("HLAinteger64Time");}

    or supply the interface that the time factory shall implement.

\end{compactstdpage}

\clearpage

% --- Standard page 309 / source PDF page 323 ---
\begin{compactstdpage}
{\ttfamily\fontsize{8.3}{9.6}\selectfont
HLAinteger64TimeFactory logicalTimeFactory =\\
LogicalTimeFactoryFactory.getLogicalTimeFactory(HLAinteger64TimeFactory.class);}

    The federate may also refer directly to the implementation provided by the RTI,

        {\ttfamily\fontsize{8.3}{9.6}\selectfont
    HLAinteger64TimeFactory logicalTimeFactory = new\\
    org.someorganization.time.HLAinteger64TimeFactoryImpl();}

    or get a list of available LogicalTimeFactories.

        {\ttfamily\fontsize{8.3}{9.6}\selectfont
    Set\textless{}LogicalTimeFactory\textgreater{} logicalTimeFactories =\\
    LogicalTimeFactoryFactory.getLogicalTimeFactories();}

    \subparagraph{Custom time types}

    An RTI implementation or a federation implementer may provide implementations of custom time types. These implementations shall be registered according to 12.11.5.24.1.

    \paragraph{Dimension upper bound and range lower and upper bounds}

    Several DDM-related services allow bounds to be gotten and set. The Get Dimension Upper Bound service gets the upper bound of a particular dimension. The Get Range Bounds and Set Range Bounds services are used to get and set the bounds of a range of a region. All bounds are represented in the Java API using a long integer. The Get Dimension Upper Bound service, as it is getting a single bound value, returns a long. The Get Range Bounds and Set Range Bounds services, as they are getting and setting a pair of values, actually return and take a RangeBounds object, which is composed of a pair of longs, lower and upper.

    \paragraph{Wall-clock time}

    The Evoke Callback and Evoke Multiple Callbacks services take arguments specifying durations of wall-clock time expressed in seconds. Wall-clock time is represented in the Java API using a double.

    \subsubsection{Memory ownership semantics}

    The following are memory management conventions for parameters to methods on RTIambassador and FederateAmbassador.

    All Java parameters, including object references, are passed by value. Therefore, there is no need to specify further conventions for primitive types.

    Unless otherwise noted, reference parameters adhere to the following convention:

    \stdliteraldashitem{The referenced object is created (or acquired) by the caller. The callee must copy during the call anything it wishes to save beyond the completion of the call.}
    \stdliteraldashitem{Unless otherwise noted, a reference returned from a method represents a new object created by the callee. The caller is free to modify the object whose reference is returned.}

    \subsection{C++}

    The C++ API is intended to achieve several goals as follows:

    \stdliteraldashitem{To provide as natural a mapping to the abstract interface specification as possible.}
\end{compactstdpage}

\clearpage

% --- Standard page 310 / source PDF page 324 ---
\begin{compactstdpage}
    \stdliteraldashitem{To prevent as many common federate programming errors as practical. This goal includes compile-time type safety as well as dynamic memory management errors.}
    \stdliteraldashitem{To provide flexibility to RTI implementers while adhering to a common API (as used by the federate programmer) and without adding significant overhead.}
    \stdliteraldashitem{To support dynamic link compatibility that allows federates to switch between RTI implementations without recompiling.}
    The following subclauses elaborate on some of these details, from the perspective of the federate programmer. This discussion is not intended as a guide for RTI implementation.

    \subsubsection{Dynamic link compatibility}

    C++ dynamic link compatibility for this specification follows the technical approach detailed in Clause 5 of SISO-STD-004.1-2004 [B5].

    The implementation of the RTI API shall be provided in a library named librti1516e followed by the extension appropriate for each platform, e.g., librti1516e.dll or librti1516e.so. A 64-bit version of the library shall have a '64' suffix, e.g., librti1516e64.so. The implementation may provide a debug version of the library with a 'd' suffix, e.g., librti1516ed.dll or librti1516e64d.dll. The librti1516e library shall contain all classes and method implementations except rti1516e::LogicalTimeFactoryFactory. The librti1516e library shall contain implementations of the standardized time classes HLAinteger64Time, HLAinteger64Interval, HLAfloat64Time, and HLAfloat64Interval, along with their factory classes HLAinteger64TimeFactory and HLAfloat64TimeFactory. The librti1516e library shall also contain an implementation of the class rti1516e::HLAlogicalTimeFactoryFactory that provides a factory for the standardized time types.

    The implementation of rti1516e::LogicalTimeFactoryFactory shall be contained in a federate-supplied library named libfedtime1516e, followed by the extension appropriate for each platform, 64-bit version, and debug version, e.g., libfedtime1516e.dll or libfedtime1516e64d.so. An RTI uses the static member function rti1516e::LogicalTimeFactoryFactory::\allowbreak{}makeLogicalTimeFactory() to create an instance of an rti1516e::LogicalTimeFactory. That factory is used in turn to create instances of federation-provided subclasses of rti1516e::LogicalTime and rti1516e::LogicalTimeInterval. The rti1516e::LogicalTimeFactoryFactory shall forward requests for other time types,    e.g.,   the standardized time types,   to rti1516e::HLAlogicalTimeFactoryFactory.

    An RTI shall provide implementations for all non-pure-virtual functions and classes listed in the RTI header files included in the appendix to this document, with the exception of rti1516e::LogicalTimeFactoryFactory.

    RTI implementations shall link directly with any required third-party libraries. This action ensures that federate developers need to link only to librti1516e and libfedtime1516e in order to use the RTI and ensures that a federate will be able to switch from one RTI implementation to another without having to re-link against the third-party libraries required by the new implementation.

    All RTIs shall support the interpretation of the FDD designator as a regular file name.

    There are several functions within the RTI API that require the use of the ostream class that is part of the C++ standard library. The legacy version of ostream, which is in the global namespace, is not supported by this API.

\end{compactstdpage}

\clearpage

% --- Standard page 311 / source PDF page 325 ---
\begin{compactstdpage}
    \paragraph{Factory mechanism}

    Implementations of the C++ API shall contain a definition of a RTIambassadorFactory class in the rti1516e namespace that contains a public createRTIambassador method. This method shall take a std::vector of std::wstring as an argument and return an std::auto\_ptr\textless{}RTIambassador\textgreater{}. The possible values for the strings in the vector argument are implementation dependent.

    \subsubsection{Memory ownership semantics}

    Many programming errors in C++ can be traced back to the improper use of the delete operator. Frequently, this tendency is a result of poor programming practice, rather than an inherent limitation of the C++ programming language. Considerable effort has been made to prevent memory corruption and leaks resulting from the misuse of the delete operator in the C++ API.

    In particular, the federate programmer never explicitly deletes objects exchanged with the RTI. Simple objects are passed by value. More complex objects are either passed by constant reference (i.e., T const \&) or passed as an auto\_ptr\textless{}T\textgreater{} object. In all three cases, the C++ compiler prevents the federate programmer from calling the delete operator on the object.

    As explained in 12.12.6, the auto\_ptr\textless{}T\textgreater{} implements the strong memory ownership and transfer semantics of the std::auto\_ptr\textless{}T\textgreater{} class\footnote{See Stroustrup [B6] for a detailed description.}.

    \subsubsection{Designators}

    The C++ interfaces that represent handles shall implement a hash() method that returns a long that is suitable for use as a hash value for use in storing handles in a hash table. The hash value may not be unique across two separate handle values but will be equal given two separate instances of the same handle value. The following are the properties of the hash value where H1 and H2 are separate instances of hash values.

    \stdliteraldashitem{H1 == H2 implies H1.hashCode() == H2.hashCode()}
    \stdliteraldashitem{H1 != H2 does not imply H1.hashCode() != H2.hashCode()}

    \subsubsection{Additional functionality}

    Implementations of the C++ API shall add the utility classes described in 12.12.4.1 and 12.12.4.2 to the rti1516e namespace.

    \paragraph{NullFederateAmbassador class}

    The NullFederateAmbassador class provides empty implementations of all methods in FederateAmbassador.

    \paragraph{Encoding helper classes}

    The encoding helper classes provide a set of classes providing helper functions for encoding datatypes according to IEEE Std 1516.2-2010.

    All encoding helper classes shall support the encode and decode methods for converting language-specific datatypes to and from standard HLA encoded byte-arrays. All noncomplex encoder classes, as listed Table 21, shall support conversion operator and operator= methods for retrieving and updating the value.


\end{compactstdpage}

\clearpage

% --- Standard page 312 / source PDF page 326 ---
\begin{compactstdpage}
    \cxiitabletwentyone

    All complex encoder classes shall support elements that make up the complex class. It shall be possible to use both complex and noncomplex encoder class instances as elements of a complex class instance.

    A method named addElement shall be supported for adding elements to an instance of such a complex class. See Table 22. Another method named get(index) for retrieving data elements shall also be supported.

    The HLAvariantRecord shall support the method addMapping to associate different encoders/decoders with different discriminants and the methods setDiscriminant and getDiscriminant to indicate which variant of the record that shall be used for composing, encoding, and decoding data elements.

    The HLAVariableArray encoder shall provide a size method that provides the number of elements.

\end{compactstdpage}

\clearpage

% --- Standard page 313 / source PDF page 327 ---
\begin{compactstdpage}
    \cxiitabletwentytwo

    \subsubsection{Services}

    In general, each service in Clause 4 through Clause 10 corresponds to one method in each API. In some cases, because of optional arguments to a service or programming language constraints, a single service may be mapped to multiple methods in an API.

    For the C++ API, the federate-initiated services invoked upon the RTI are members of the RTIambassador class, and RTI-initiated service callbacks on the federate are members of the FederateAmbassador class. The RTI and federate ambassador services in the C++ API that do not follow a one-to-one mapping with the services in Clause 4 through Clause 10 of the specification are described in 12.12.5.1 through 12.12.5.34.

    \paragraph{Create Federation Execution section}

    The set of FOM module designators argument is implemented as a std::vector of std::wstrings where each item in the vector is a full path to a file containing one FOM module. The optional MIM designator is implemented as an argument with a default value of the empty string.

    In addition to the supplied arguments listed with the description of this service, the C++ API requires the federate to supply a wstring logicalTimeImplementationName that contains the name of the logical time implementation. In addition to the specified exceptions listed with the description of this service, the C++ API may also throw a CouldNotCreateLogicalTimeFactory exception. This exception indicates that the RTI was not able to create an instance of the logical time factory class that was named in the wstring logicalTimeImplementationName argument.

    \paragraph{Join Federation Execution section}

    The optional set of FOM module designators argument is implemented as a std::vector of std::wstrings where each item in the vector is a full path to a file containing one FOM module. If the federate does not need to load additional FOM modules, then an empty std::vector should be passed to the joinFederationExecution method.

    In addition to the specified exceptions listed with the description of this service, the C++ API may also throw a CouldNotCreateLogicalTimeFactory or CallNotAllowedFromWithinCallback exception. The CouldNotCreateLogicalTimeFactory exception indicates that the RTI was not able to create an instance of the logical time factory class that was named in the wstring logicalTimeImplementationName argument to the Create Federation Execution service. The CallNotAllowedFromWithinCallback exception indicates that this RTI ambassador service call cannot be made while the federate is processing a federate ambassador callback.

\end{compactstdpage}

\clearpage

% --- Standard page 314 / source PDF page 328 ---
\begin{compactstdpage}
    \paragraph{Register Federation Synchronization Point service}

    The optional argument representing the set of joined federate designators is implemented in the C++ API as two registerFederationSynchronizationPoint() methods: one that takes an FederateHandleSet as an argument, and one that does not.

    \paragraph{Confirm Synchronization Point Registration \rtiinitiated{} service}

    The registration-success indicator is implemented in the C++ API as two distinct methods: synchronizationPointRegistrationSucceeded()               and synchronizationPointRegistrationFailed(). A synchronization failure reason (discussed in 12.12.6.10) is passed as an argument to synchronizationPointRegistrationFailed().

    \paragraph{Synchronization Point Achieved service}

    The optional synchronization-success indicator parameter is a required argument in the C++ API.

    \paragraph{Request Federation Save service}

    The optional timestamp argument is implemented as two methods, one of which takes an LogicalTime as an argument.

    \paragraph{Initiate Federate Save \rtiinitiated{} service}

    The optional timestamp argument is implemented as two methods, one of which takes an LogicalTime as an argument.

    \paragraph{Federate Save Complete service}

    This service corresponds to two methods in the C++ API: federateSaveComplete()                                                    and federateSaveNotComplete().

    \paragraph{Federation Saved \rtiinitiated{} service}

    This service corresponds to two methods in the C++ API: federationSaved() and federationNotSaved(). A Save Failure Reason (discussed in 12.12.6.11) is passed as an argument to federationNotSaved().

    \paragraph{Confirm Federation Restoration Request \rtiinitiated{} service}

    This service is mapped to the following two methods in the C++ API: requestFederationRestoreSucceeded() and requestFederationRestoreFailed().

    \paragraph{Federate Restore Complete service}

    This service is mapped to the following two methods in the C++ API: federateRestoreComplete() and federateRestoreNotComplete().

    \paragraph{Federation Restored \rtiinitiated{} service}

    This service corresponds to two methods in the C++ API: federationRestored() and federationNotRestored(). A Restore Failure Reason (discussed in 12.12.6.12) is passed as an argument to federationNotRestored().

\end{compactstdpage}

\clearpage

% --- Standard page 315 / source PDF page 329 ---
\begin{compactstdpage}
    \paragraph{Connect service}

    In addition to the supplied arguments listed with the description of this service, the C++ API requires the federate to supply an instance of the interface FederateAmbassador. The optional localSettingsDesignator parameter is required and can contain RTI implementation-specific configuration information. An empty string for the localSettingsDesignator parameter shall be used to indicate the default local settings.

    In addition to the specified exceptions listed with the description of this service, the C++ API may also throw a CallNotAllowedFromWithinCallback exception. The CallNotAllowedFromWithinCallback exception indicates that this RTI ambassador service call cannot be made while the federate is processing a federate ambassador callback.

    \paragraph{Unpublish Object Class Attributes service}

    The optional argument representing the set of attribute designators is implemented in the C++ API as two methods: unpublishObjectClass() and unpublishObjectClassAttributes(). The latter of the two takes an AttributeHandleSet as an argument.

    \paragraph{Unsubscribe Object Class Attributes service}

    The optional argument representing the set of attribute designators is implemented in the C++ API as two methods: unsubscribeObjectClass() and unsubscribeObjectClassAttributes(). The latter of the two takes an AttributeHandleSet as an argument.

    \paragraph{Reserve Object Instance Name service}

    The optional single name parameter or list of names parameter is implemented in the C++API as two methods: reserveObjectInstanceName() and reserveMultipleObjectInstanceName().

    \paragraph{Object Instance Name Reserved \rtiinitiated{} service}

    This service corresponds to two methods in the C++ API: objectInstanceNameReservationSucceeded() and objectInstanceNameReservationFailed().

    \paragraph{Multiple Object Instance Name Reserved \rtiinitiated{} service}

    This service corresponds to two methods in the C++ API: multipleObjectInstanceNameReservationSucceeded() and multipleObjectInstanceNameReservationFailed().

    \paragraph{Register Object Instance service}

    The optional argument representing the object instance name is implemented in the C++ API as two registerObjectInstance() methods: one that takes a wstring as an argument for the object instance name, and one that does not.

    \paragraph{Discover Object Instance service}

    The optional argument representing the producing joined federate designator is implemented in the C++ API as two discoverObjectInstance() methods: one that takes a FederateHandle as an argument for the producing joined federate designator, and one that does not.

\end{compactstdpage}

\clearpage

% --- Standard page 316 / source PDF page 330 ---
\begin{compactstdpage}
    \paragraph{Update Attribute Values service}

    The optional argument representing the timestamp is implemented in the C++ API as two uptdateAttributeValues() methods: one that takes a LogicalTime as an argument for the timestamp, and one that does not. The version that takes a LogicalTime also returns a MessageRetractionHandle for the message retraction designator.

    \paragraph{Reflect Attribute Values \rtiinitiated{} service}

    This service corresponds to three overloaded versions of the reflectAttributeValues() method in the C++ API. Which version the RTI calls is dependent on if the updateAttributeValues() was called with a timestamp and if the message was received TSO. See 12.8 for details how to determine which version the RTI will call.

    \paragraph{Send Interaction service}

    The optional argument representing the timestamp is implemented in the C++ API as two sendInteraction() methods: one that takes a LogicalTime as an argument for the timestamp, and one that does not. The version that takes a LogicalTime also returns a MessageRetractionHandle for the message retraction designator.

    \paragraph{Receive Interaction \rtiinitiated{} service}

    This service corresponds to three overloaded versions of the receiveInteraction() method in the C++ API. Which version the RTI calls is dependent on if the sendInteraction() was called with a timestamp and if the message was received TSO. See 12.9 for details how to determine which version the RTI will call.

    \paragraph{Delete Object Instance service}

    The optional argument representing the timestamp is implemented in the C++ API as two deleteObjectInstance() methods: one that takes a LogicalTime as an argument for the timestamp, and one that does not. The version that takes a LogicalTime also returns a MessageRetractionHandle for the message retraction designator.

    \paragraph{Remove Object Instance \rtiinitiated{} service}

    This service corresponds to three overloaded versions of the removeObjectInstance() method in the C++ API. See 12.10 for details on how to determine which version the RTI will call.

    \paragraph{Request Attribute Value Update service}

    This service corresponds to two overloaded methods named requestAttributeValueUpdate() in the C++ API. One takes an ObjectInstanceHandle to represent the object instance designator. The other takes an ObjectClassHandle to represent the object class designator.

    \paragraph{Inform Attribute Ownership \rtiinitiated{} service}

    In the C++ API, this service corresponds to three distinct methods: attributeIsNotOwned(), attributeIsOwnedByRTI(), and informAttributeOwnership(). All three methods take an ObjectInstanceHandle to represent the object instance designator and an AttributeHandle to represent the attribute designator as arguments. In addition, informAttributeOwnership() takes a FederateHandle as an argument to represent the federate designator.

\end{compactstdpage}

\clearpage

% --- Standard page 317 / source PDF page 331 ---
\begin{compactstdpage}
    \paragraph{Register Object Instance With Regions service}

    This service corresponds to two overloaded versions of the registerObjectInstanceWithRegions() method in the C++ API. The second version takes a wstring as an argument to represent the optional object instance name.

    \paragraph{Send Interaction With Regions service}

    The optional argument representing the timestamp is implemented in the C++ API as two sendInteractionWithRegions() methods: one that takes a LogicalTime as an argument for the timestamp, and one that does not. The version that takes a LogicalTime also returns a MessageRetractionHandle for the message retraction designator.

    \paragraph{Resign Federation Execution service}

    In addition to the specified exceptions listed with the description of this service, the C++ API may also throw a CallNotAllowedFromWithinCallback exception. This exception indicates that this RTI ambassador service call cannot be made while the federate is processing a federate ambassador callback.

    \paragraph{Disconnect service}

    In addition to the specified exceptions listed with the description of this service, the C++ API may also throw a CallNotAllowedFromWithinCallback exception. This exception indicates that this RTI ambassador service call cannot be made while the federate is processing a federate ambassador callback.

    \paragraph{Evoke Callback service}

    In addition to the specified exceptions listed with the description of this service, the C++ API may also throw a CallNotAllowedFromWithinCallback exception. This exception indicates that this RTI ambassador service call cannot be made while the federate is processing a federate ambassador callback.

    \paragraph{Evoke Multiple Callbacks service}

    In addition to the specified exceptions listed with the description of this service, the C++ API may also throw a CallNotAllowedFromWithinCallback exception. This exception indicates that this RTI ambassador service call cannot be made while the federate is processing a federate ambassador callback.

    \subsubsection{Abstract datatypes}

    Subclause 1.4.2 presents general conventions used to describe the abstract datatypes used in Clause 4 through Clause 10. In 12.12.6.1 through 12.12.6.26, the abstract datatypes are further detailed by mapping an individual datatype to its implementation in the C++ API. Note that the datatypes from Clause 11 are explained in that clause and are not repeated here. Additionally, designators are a subject in their own right and are covered in 12.2.

    The ISO/IEC committee ratified the C++ standard as ISO/IEC 14882 [B3]. The HLA C++ API utilizes standard C++ constructs to the fullest extent possible. For example, std::wstring is used to represent character strings that might previously have been represented with the C-style construct, pointer to char (i.e., char *). Additionally, several standard C++ containers are used to represent the sets, constrained sets of pairs, and collections described in Clause 3. These containers are std::set\textless{}T\textgreater{}, std::map\textless{}T\textgreater{}, and std::vector\textless{}T\textgreater{}. Also, the C++ standard std::auto\_ptr\textless{}T\textgreater{} is used to implement strict memory ownership transfer when passing instances of dynamically allocated datatypes across API.

\end{compactstdpage}

\clearpage

% --- Standard page 318 / source PDF page 332 ---
\begin{compactstdpage}
    A description of the nature and properties of these standard C++ constructs is well beyond the scope of this clause. Instead, readers are referred to the C++ ISO standard as well as any of the vast assortment of C++ references. Two such excellent references are Austern [B1] and Stroustrup [B6].

    \paragraph{Names, labels, and strings}

    The names, labels, and strings described in Clause 4 through Clause 10 of this specification are mapped to the wstring type.

    \paragraph{Boolean values}

    Standard C++ provides a fundamental type, bool, intended for the representation of boolean values. This fundamental type is used for all boolean values in the C++ interface.

    \paragraph{Exceptions}

    In keeping with the philosophy outlined in 12.12.6, exceptions in the C++ API have an interface identical with std::exception. In fact, all exceptions in the C++ API inherit from a class named Exception. This hierarchy allows all exceptions in the C++ API to be caught as a single base class.

    \paragraph{Handles}

    Most designators in the C++ API are handles. Each handle is a distinct type, but they all share a common interface, defined by the DEFINE\_HANDLE\_CLASS macro and specialized into one of the following concrete classes: FederateHandle, ObjectClassHandle, InteractionClassHandle, ObjectInstanceHandle, AttributeHandle, ParameterHandle, DimensionHandle, MessageRetractionHandle, and RegionHandle.

    Two instances of each class with identical types can be compared for equality or order (i.e., less than), and one can be assigned to the other. A copy constructor is provided as well. These handles are always generated by an RTI implementation; however, a default constructor is provided for convenience. Additionally, an encode() method is provided that returns an encoded form of a handle, VariableLengthData.

    VariableLengthData has a data() method that returns an untyped pointer (i.e., a void const *) as well as a size() method that returns the size in bytes (represented as a unsigned long) of the data returned by the data() method. A copy constructor and a copy-assignment operator are also provided.

    Using this interface, it is possible for the RTI user to save and restore any handle provided by the RTI as well as to pass handles as attribute values or parameter values among federates. To send a handle as an attribute value or as a parameter value, the encode() method of the handle should be invoked by the federate programmer. The returned VariableLengthData class has a data() method and a size() method that can be used to fill in the attribute value or parameter value. Once delivered, the encoded handle can be reconstituted into an instance of its corresponding handle class using that handle class's constructor.

    Note well that the pointer value returned by the data() method is not required to have any particular byte alignment. Consequently, naively attempting to simply type cast this pointer and assign its contents to some type may generate an exception in the processor's memory management unit and likely result in the crash of the program. If for some reason it is desirable to perform this type of operation, a function such as memcpy() should be used.

    The VariableLengthData class has a constructor that requires the user to provide an untyped pointer (i.e., a void const *) as well as the size in bytes (represented as a unsigned long) of the data to which it points. Methods to return these values are provided as well as is a copy constructor and a copy-assignment operator.

\end{compactstdpage}

\clearpage

% --- Standard page 319 / source PDF page 333 ---
\begin{compactstdpage}
    Handle classes also provide a toString() method that returns a printable form of a handle. The string produced, however, is not guaranteed to be the same as the name to which the handle corresponds (assuming there is a name that corresponds to that kind of handle). If the user needs the name form of a handle, employ the appropriate support service from Clause 10.

    \paragraph{Resignation directive}

    The Resign Federation Execution service allows the specification of one of six directives. To provide this specification, a ResignAction enumeration is defined. The six different directives of Resign Federation Execution are defined as the six enumerated values that instances of the ResignAction enumeration may take. The ResignAction enumeration has the following values:

    \stdliteraldashitem{UNCONDITIONALLY\_DIVEST\_ATTRIBUTES}
    \stdliteraldashitem{DELETE\_OBJECTS}
    \stdliteraldashitem{CANCEL\_PENDING\_OWNERSHIP\_ACQUISITIONS}
    \stdliteraldashitem{DELETE\_OBJECTS\_THEN\_DIVEST}
    \stdliteraldashitem{CANCEL\_THEN\_DELETE\_THEN\_DIVEST}
    \stdliteraldashitem{NO\_ACTION}

    \paragraph{User-supplied tag}

    Several services in the HLA allow the user to supply additional information to be transported by the RTI in an opaque manner, such as the Update Attribute Values service. The C++ API uses the VariableLengthData class to achieve this goal.

    \paragraph{Sets of designators}

    Sets of designators are implemented using the standard C++ set container with the appropriate handle type. For convenience, each different set of designators has been aliased using a typedef. To ensure maximum portability between RTI implementations, federate programmers should use only this typedef. The list of set types with their corresponding typedefs are as follows:

    \stdliteraldashitem{A set of federate designators uses std::set\textless{}FederateHandle\textgreater{}, and the typedef FederateHandleSet should be used.}
    \stdliteraldashitem{A set of attribute designators uses std::set\textless{}AttributeHandle\textgreater{}, and the typedef AttributeHandleSet should be used.}
    \stdliteraldashitem{A set of attribute designators uses std::set\textless{}ParameterHandle\textgreater{}, and the typedef ParameterHandleSet should be used.}
    \stdliteraldashitem{A set of dimension designators uses std::set\textless{}DimensionHandle\textgreater{}, and the typedef DimensionHandleSet should be used.}
    \stdliteraldashitem{A set of region designators uses std::set\textless{}RegionHandle\textgreater{}, and the typedef RegionHandleSet should be used.}

    \paragraph{Success indicators}

    Several of the requests a federate may make of an RTI shall be acknowledged by the RTI to the federate. This acknowledgment can indicate that the request succeeded or failed. In the C++ API, each request acknowledgment corresponds to a positive and a negative callback method on the FederateAmbassador,             e.g.,  federationSaved()           and federationNotSaved(), requestFederationRestoreSucceeded() and requestFederationRestoreFailed(), and federationRestored() and federationNotRestored().

\end{compactstdpage}

\clearpage

% --- Standard page 320 / source PDF page 334 ---
\begin{compactstdpage}
    \paragraph{Callback model}

    When the connect() method is invoked by a federate, it must supply an instance of the CallbackModel enumeration as an argument. The CallbackModel enumeration provides the following values:

    \stdliteraldashitem{HLA\_IMMEDIATE}
    \stdliteraldashitem{HLA\_EVOKED}

    \paragraph{Synchronization failure reason}

    When the synchronizationPointRegistrationFailed() method is invoked by the RTI on a federate, the RTI provides an instance of a SynchronizationFailureReason enumeration as an argument. The SynchronizationFailureReason enumeration provides the following values:

    \stdliteraldashitem{SYNCHRONIZATION\_POINT\_LABEL\_NOT\_UNIQUE}
    \stdliteraldashitem{SYNCHRONIZATION\_SET\_MEMBER\_NOT\_JOINED}

    \paragraph{Save Failure Reason}

    When the federationNotSaved() method is invoked by the RTI on a federate, the RTI provides an instance of a SaveFailureReason enumeration as an argument. The SaveFailureReason enumeration provides the following values:

    \stdliteraldashitem{UNABLE\_TO\_SAVE}
    \stdliteraldashitem{FEDERATE\_REPORTED\_FAILURE\_DURING\_SAVE}
    \stdliteraldashitem{FEDERATE\_RESIGNED\_DURING\_SAVE}
    \stdliteraldashitem{RTI\_DETECTED\_FAILURE\_DURING\_SAVE}
    \stdliteraldashitem{SAVE\_TIME\_CANNOT\_BE\_HONORED}
    \stdliteraldashitem{SAVE\_ABORTED}

    \paragraph{Restore Failure Reason}

    When the federationNotRestored() method is invoked by the RTI on a federate, the RTI provides an instance of a RestoreFailureReason enumeration as an argument. The RestoreFailureReason enumeration provides the following values:

    \stdliteraldashitem{RTI\_UNABLE\_TO\_RESTORE}
    \stdliteraldashitem{FEDERATE\_REPORTED\_FAILURE\_DURING\_RESTORE}
    \stdliteraldashitem{FEDERATE\_RESIGNED\_DURING\_RESTORE}
    \stdliteraldashitem{RTI\_DETECTED\_FAILURE\_DURING\_RESTORE}
    \stdliteraldashitem{RESTORE\_ABORTED}

    \paragraph{List of joined federates save status}

    The list of joined federates and the save status of each is implemented as a vector of federate designator and save status pairs. In the C++ API, this corresponds to an vector\textless{}pair\textless{} FederateHandle, SaveStatus\textgreater{}\textgreater{}. For convenience, this type has been aliased using a typedef to the name FederateHandleSaveStatusPairVector. To ensure maximum portability between RTI implementations, federate programmers should use only this typedef.

    SaveStatus is an enumeration that provides the following values:

    \stdliteraldashitem{NO\_SAVE\_IN\_PROGRESS}
\end{compactstdpage}

\clearpage

% --- Standard page 321 / source PDF page 335 ---
\begin{compactstdpage}
    \stdliteraldashitem{FEDERATE\_INSTRUCTED\_TO\_SAVE}
    \stdliteraldashitem{FEDERATE\_SAVING}
    \stdliteraldashitem{FEDERATE\_WAITING\_FOR\_FEDERATION\_TO\_SAVE}

    \paragraph{List of joined federates restore status}

    The list of joined federates and the restore status of each is implemented as a vector of federate designator and restore status pairs. In the C++ API, this corresponds to an vector\textless{}pair\textless{} FederateHandle, RestoreStatus\textgreater{}\textgreater{}. For convenience, this type has been aliased using a typedef to the name Federate HandleRestoreStatusPairVector. To ensure maximum portability between RTI implementations, federate programmers should use only this typedef.

    RestoreStatus is an enumeration that provides the following values:

    \stdliteraldashitem{NO\_RESTORE\_IN\_PROGRESS}
    \stdliteraldashitem{FEDERATE\_RESTORE\_REQUEST\_PENDING}
    \stdliteraldashitem{FEDERATE\_WAITING\_FOR\_RESTORE\_TO\_BEGIN}
    \stdliteraldashitem{FEDERATE\_PREPARED\_TO\_RESTORE}
    \stdliteraldashitem{FEDERATE\_RESTORING}
    \stdliteraldashitem{FEDERATE\_WAITING\_FOR\_FEDERATION\_TO\_RESTORE}

    \paragraph{Passive subscription indicator}

    The Subscribe Object Class Attributes, Subscribe Interaction Class, Subscribe Object Class Attributes With Regions, and Subscribe Interaction Class With Regions services each may optionally provide an bool value to indicate whether the subscription is active (true means the subscription is active). In the C++ API, this is a default parameter to the methods corresponding to each of the services. The default value is true.

    \paragraph{Constrained set of attribute designator and value pairs}

    The Update Attribute Values and Reflect Attribute Values \rtiinitiated{} services require a constrained set of attribute designator and value pairs as arguments. A constrained set of attribute designator and value pairs is implemented in the C++ API as an map\textless{}AttributeHandle, VariableLengthData\textgreater{}. For convenience, this type has been aliased using a typedef to the name AttributeHandleValueMap. To ensure maximum portability between RTI implementations, federate programmers should use only this typedef.

    \paragraph{Constrained set of interaction parameter designator and value pairs}

    The Send Interaction, Receive Interaction \rtiinitiated{}, and Send Interaction With Regions services require a constrained set of parameter designator and value pairs as arguments. A constrained set of parameter designator and value pairs is implemented in the C++ API as an map\textless{}ParameterHandle, VariableLengthData\textgreater{}. For convenience, this type has been aliased using a typedef to the name ParameterHandleValueMap. To ensure maximum portability between RTI implementations, federate programmers should use only this typedef.

    \paragraph{Message order type}

    Message order types are represented an enumeration named OrderType, which has the following values: RECEIVE and TIMESTAMP. These values can be passed to the RTI or compared to an instance of an OrderType object returned by the RTI.

\end{compactstdpage}

\clearpage

% --- Standard page 322 / source PDF page 336 ---
\begin{compactstdpage}
    \paragraph{Transportation type}

    As with message order types described above, transportation types are represented by an enumeration. An RTI shall provide the following enumerations: RTI\_RELIABLE and RTI\_BEST\_EFFORT. These enumerations may be passed to the RTI or compared to an instance of a TransportationType enumeration returned by the RTI.

    It is possible that a FDD may contain additional transportation types. In this case, the RTI implementation shall extend the TransportationType enumerated values to provide a unique value for each new transportation type. These new transportation type enumerated values shall be retrievable from the RTI via the getTransportationType service. In no case shall the existing TransportationType enumerations or the Enums.h header file be modified or eliminated.

    \paragraph{Ownership designator}

    The Inform Attribute Ownership \rtiinitiated{} service supplies an ownership designator as one of its arguments. In the C++ API, this corresponds to three distinct methods: attributeIsNotOwned(), attributeIsOwnedByRTI(), and informAttributeOwnership(). All three methods take an ObjectInstanceHandle and an AttributeHandle as arguments. In addition, informAttributeOwnership() takes an FederateHandle as an argument. Together, these three methods combine to represent the ownership designator of the Inform Attribute Ownership \rtiinitiated{} service.

    \paragraph{Definition indicator}

    The Query GALT and Query LITS services are required to return an indicator to define whether their return value is valid. Since both services return values of type auto\_ptr\textless{}LogicalTime\textgreater{}, the definition indicator is trivial: if the returned value is NULL, the value is undefined (invalid); otherwise, the returned value is defined (valid).

    \paragraph{Optional message retraction designator}

    The Update Attribute Values, Send Interaction, Send Interaction with Regions, and Delete Object Instance services may return a message retraction designator under certain circumstances. In the C++ API, this is implemented much like the definition indicator described above. The methods corresponding to these services always return an auto\_ptr\textless{}\allowbreak MessageRetractionHandle\textgreater{}. If the returned value is NULL, then the operation did not generate a message retraction designator; otherwise, the auto\_ptr\textless{}\allowbreak MessageRetractionHandle\textgreater{} holds the returned message retraction designator.

    \paragraph{Collection of attribute designator set and region designator set pairs}

    Several services require the federate to provide the RTI with a collection of attribute designator set and region designator set pairs. In the C++ API, this corresponds to a vector\textless{} pair\textless{}AttributeHandleSet, RegionHandleSet\textgreater{}\textgreater{}. For convenience, this type has been aliased using a typedef to the name AttributeHandleSetRegionHandleSetPair Vector. To ensure maximum portability between RTI implementations, federate programmers should use only this typedef.

    \paragraph{Logical time, timestamps, and lookahead}

    The C++ API defines the logical time-related classes as abstract base classes. The method for obtaining a printable string representation of a LogicalTime and a LogicalTimeInterval is toString() and it returns a wstring. The encode() and decode()methods defined for the LogicalTime and LogicalTimeInterval classes use objects of type VariableLengthDataType for the variable array. The C++ API also provides encode() and decode()methods that allow encoding and decoding

\end{compactstdpage}

\clearpage

% --- Standard page 323 / source PDF page 337 ---
\begin{compactstdpage}
    directly into a void * buffer. The encoding and decoding methods are available on the LogicalTime, LogicalTimeInterval, and LogicalTimeFactory classes.

    In addition to the methods mentioned in 12.3, the C++ LogicalTime interface provides the following methods:

    \stdliteraldashitem{setInitial() that sets the current instance to the initial value}
    \stdliteraldashitem{setFinal() that sets the current instance to the final value}
    \stdliteraldashitem{encodedLength() that returns the number of bytes needed to encode the current instance}
    \stdliteraldashitem{implementationName() that returns a std::wstring with the name of the logical time implementation that is in use by the createFederationExecution method}
    In addition to the methods mentioned in 12.3, the C++ LogicalTimeInterval interface provides the following methods:

    \stdliteraldashitem{setZero() that sets the current instance to the zero value}
    \stdliteraldashitem{setEpsilon() that sets the current instance to the epsilon value}
    \stdliteraldashitem{encodedLength() that returns the number of bytes needed to encode the current instance}
    \stdliteraldashitem{implementationName() that returns a std::wstring with the name of the logical time implementation that is in use by the createFederationExecution method}
    In addition to the methods mentioned in 12.3, the C++ LogicalTimeFactory interface provides the following methods:

    \stdliteraldashitem{setZero() that sets the current instance to the zero value}
    \stdliteraldashitem{makeInitial() to create an instance of a LogicalTime with the initial value}
    \stdliteraldashitem{makeFinal() to create an instance of a LogicalTime with the final value}
    \stdliteraldashitem{makeZero() to create an instance of a LogicalTimeInterval with the value of zero}
    \stdliteraldashitem{makeEpsilon() to create an instance of a LogicalTimeInterval with the value of epsilon}
    A C++ library that implements the logical time classes shall supply a exported static method with the following signature that allows an RTI implementation to create an instance of the LogicalTimeFactory(). Federates should not use this interface to obtain an LogicalTimeFactory but should instead use the getTimeFactory() method on the RTI ambassador.

        {\ttfamily\fontsize{8.3}{9.6}\selectfont
    static std::auto\_ptr\textless{}LogicalTimeFactory\textgreater{}\\
    makeLogicalTimeFactory(std::wstring const \& implementationName);}

    \subparagraph{Standardized time types}

    The implementations of the standardized time types shall be placed in the hla.rti1516e namespace and have the names specified in Table 23.

    \cxiitabletwentythree

\end{compactstdpage}

\clearpage

% --- Standard page 324 / source PDF page 338 ---
\begin{compactstdpage}
    \paragraph{Dimension upper bound and range lower and upper bounds}

    Several DDM-related services allow bounds to be gotten and set. The Get Dimension Upper Bound service gets the upper bound of a particular dimension. The remaining services are used to get and set the bounds of ranges of a region. All bounds are represented in the C++ API using an unsigned integer.

    Range lower and upper bounds are gathered into a RangeBounds object.

    \paragraph{Wall-clock time}

    The Evoke Callback and Evoke Multiple Callbacks services take arguments specifying durations of wall-clock time expressed in seconds. Wall-clock time is represented in the C++ API using a double.

    \subsection{Web Service Definition Language (WSDL)}

    \subsubsection{Overview}

    \paragraph{Standards and terminology}

    The Web Service API is specified using the WSDL version 1.1. WSDL uses XML Schema and is based on the SOAP 1.1 specification that in turn uses XML version 1.0. The design of the Web Service API follows the guidelines of the Web Service Interoperability (WS-I) standard version 1.1. Full references to the above standards are shown in Table 24.

    \cxiitabletwentyfour

    The following terms in the text are used based on their WSDL and Web Service sense unless otherwise indicated: service, consumer, provider, nillable, port, transportation and http.

    The namespace http://sisostds.org/hla/evolved/wsdl/v1 is used for this WSDL.

    \paragraph{HLA Web service consumers and providers}

    This API differs from the C++ and Java APIs since it is not actually a programming API but a protocol used for communicating between a federate and the RTI. The federate in this context is considered to be a Web Service Consumer. The Web Service is provided by a software component called the Web Service Provider RTI Component (WSPRC), often referred to as the Provider.

\end{compactstdpage}

\clearpage

% --- Standard page 325 / source PDF page 339 ---
\begin{compactstdpage}
    Federates that wish to interoperate must connect to RTI Components [Local Runtime Components (LRCs) or WSPRCs] that are able to establish a common federation execution. No additional restrictions are imposed on the topology or implementation of the RTI. One RTI may provide several WSPRCs. One WSPRC may support several simultaneously connected federates. One WSPRC may provide several ports that Consumers can connect to.

    \paragraph{Implementation language considerations}

    The Web Service API covers the full HLA interface specification. A Web Service federate or RTI may be implemented in any language and environment capable of Web Service communication.

    There are several ways to implement a federate using the WSDL API. It is, for example, possible to generate code for several languages using WSDL-to-source code generators. Note that if code is generated for C++ and Java, this code is likely to differ from the corresponding standard HLA APIs.

    \paragraph{Transportation protocols and ports}

    The HLA Web Service API shall make use of Hypertext Transfer Protocol (HTTP) for transportation.

    The WSDL API in Annex D uses a placeholder for the port, which is the connection point for a Consumer. When implementing the WSDL API, this needs to be replaced with the port where the HLA Web Service Provider is available.

    \subsubsection{Services}

    \paragraph{Optional parameters and polymorphism}

    Parameters that are described as optional in this standard are defined as Nillable in the Web Service API.

    In case of polymorphism where the same HLA service can be called with parameters of different types, these HLA services are divided into several Web services.

    \paragraph{Exceptions}

    All exceptions are implemented as SOAP Fault elements and are handled according to the SOAP standard see 12.13.1.1). One SOAP Fault element is generated for each exception.

    \paragraph{FDD and connection parameters}

    When invoking the Create Federation Execution service, the FDD shall be referenced by a set of URLs. The URLs in this set shall point to the FOM modules that apply to that FDD. These FOM modules shall be accessible for the WSPRC.

    The local settings designator in the connect call shall be used to provide information to the WSPRC to be used for connecting to other RTI components.

    \paragraph{Modified services}

    The Request Attribute Value Update service may take either an object instance or a class designator as its first parameter. This service is implemented as two different services in the WSDL API:

    \stdliteraldashitem{Request Object Instance Attribute Value Update}
    \stdliteraldashitem{Request Object Class Attribute Value Update}
\end{compactstdpage}

\clearpage

% --- Standard page 326 / source PDF page 340 ---
\begin{compactstdpage}
    The Inform Attribute Ownership service (callback) is implemented using the following calls:

    \stdliteraldashitem{Inform Attribute Ownership}
    \stdliteraldashitem{Attribute Is Not Owned}
    \stdliteraldashitem{Attribute Owned By RTI}
    The Federate Save Complete service is implemented using the following calls:

    \stdliteraldashitem{Federate Save Complete}
    \stdliteraldashitem{Federate Save Not Complete}
    The Federate Restore Complete service is implemented using the following calls:

    \stdliteraldashitem{Federate Restore Complete}
    \stdliteraldashitem{Federate Restore Not Complete}

    \subsubsection{Datatypes}

    \paragraph{Encoding of parameters in the interface specification}

    Table 25 summarizes how parameters shall be represented in the Web Service API.

    The xsd prefix indicates the http://www.w3.org/2001/XMLSchema namespace.

    \cxiitabletwentyfivea

\end{compactstdpage}

\clearpage

% --- Standard page 327 / source PDF page 341 ---
\begin{compactstdpage}
    \cxiitabletwentyfiveb

\end{compactstdpage}

\clearpage

% --- Standard page 328 / source PDF page 342 ---
\begin{compactstdpage}
    \cxiitabletwentyfivec

    \paragraph{Encoding of FDD-defined data values}

    Certain data shall first be encoded according to the HLA OMT and then encoded using xsd:base64Binary before being sent. This requirement applies to the following data:

    \stdliteraldashitem{Attribute values}
    \stdliteraldashitem{Parameter values}
    \stdliteraldashitem{Logical timestamps}
\end{compactstdpage}

\clearpage

% --- Standard page 329 / source PDF page 343 ---
\begin{compactstdpage}
    Web Service federates have the option of using the encoding helpers available in the Java and C++ APIs for the encoding. In addition to this option, a time-regulating and/or time-constrained federate needs access to a time class implementation that is compatible with the one used in the federation.

    \paragraph{Time representation and implementation}

    The WSDL API requires no standardized local RTI component for a federate; in other words, no local time or interval classes are provided. The time values are exchanged between the federate and the RTI as byte arrays. An RTI shall support federates using the WSDL API that use the standardized support services below for encoding, decoding, and updating timestamps as well as federates performing the corresponding operations without using these services.

    Implementers of the WSDL API shall provide, as part of their implementation, implementations of encoding and decoding services for two standardized logical time types: HLAinteger64Time and HLAfloat64Time using the service names shown in Table 26.

    \cxiitabletwentysix

    \subsubsection{Handles}

    \paragraph{Handles for FDD data}

    Unlike the Java and C++ APIs, the WSDL API does not use handles for FDD data. Instead, the WSDL API uses fully qualified names as designated in Table 27.

    \cxiitabletwentysevena

\end{compactstdpage}

\clearpage

% --- Standard page 330 / source PDF page 344 ---
\begin{compactstdpage}
    \cxiitabletwentysevenb

    To simplify porting between the WSDL API and other APIs, the support services for FDD data handles are still available. Services that return handles in the Java or C++ API shall instead return the fully qualified name in the WSDL API. Services that take handles as input in the Java or C++ API shall instead take the fully qualified name as input in the WSDL API. Exceptions for incorrect values shall be thrown as for the C++ and Java API. The following services take a handle or fully qualified name, respectively:

    \stdliteraldashitem{Get Object Class Handle / Get Object Class Name}
    \stdliteraldashitem{Get Attribute Handle / Get Attribute Name}
    \stdliteraldashitem{Get Interaction Class Handle / Get Interaction Class Name}
    \stdliteraldashitem{Get Parameter Handle / Get Parameter Name}
    \stdliteraldashitem{Get Dimension Handle / Get Dimension Name}
    \stdliteraldashitem{Get Object Instance Handle / Get Object Instance Name}
    \stdliteraldashitem{Get Federate Handle / Get Federate Name}

    \paragraph{Services for encoding and decoding handles}

    The WSDL API still offers the ability to exchange handles with other federates. The principle is the same as for the C++ and Java API: a handle must be encoded prior to sending and decoded by the receiver. The services shown in Table 28 are provided to support this activity.

    \cxiitabletwentyeight

\end{compactstdpage}

\clearpage

% --- Standard page 331 / source PDF page 345 ---
\begin{compactstdpage}
    \subsubsection{Special Web Service considerations}

    \paragraph{Delivery of callbacks}

    The current WSDL API supports only requests from the Consumer (federate) and responses from the Provider (WSPRC). To get callbacks from the WSPRC to the Consumer, the following service needs to be called:

    \stdliteraldashitem{EvokeMultipleCallbacks (Maximum Callback Count).}
    The return value of this service is a CallbackArray element containing a \textless{}xsd:sequence\textgreater{} of callback elements as specified in the WSDL. A maximum callback count is supplied. The minimum and maximum time is removed. The RTI shall return no more than the specified number of callbacks.

    The pending indicator is returned in both cases to provide the number of pending callbacks.

    Note: The Evoke Callback function is unavailable in the WSDL API since the call EvokeMultipleCallbacks does the same function.

    How frequently these services need to be called may vary between different federates, federations, and scenarios. The Consumer may also need to take the session timeout (described below) into consideration.

    \paragraph{Session management}

    An ongoing connection between a federate and the RTI in the C++ and Java API is constituted by a federate ambassador and RTI ambassador pair. In the WSDL API, an ongoing connection is instead implemented through a session identifier (ID). The session ID is initially provided by the WSPRC when the Consumer issues the connect call unless it already has a valid session ID. This session ID is subsequently used to identify this ongoing connection between the WSPRC and the Consumer. The session ID is sent as a SOAP-header in requests and responses.

    The WSPRC (RTI) may terminate the session and dispose of the session data if the Consumer has not called the WSPRC for a certain time. The session is then considered unavailable. This situation may, for example, happen if the Consumer has terminated unexpectedly.

    There is no prescribed value for this timeout value; in other words, different schemas may be implemented including an infinite timeout value.

    A session can also become unavailable for other reasons, for example, if the WSPRC fails and is restarted.

    \paragraph{Fault tolerance}

    Fault tolerance considerations must take into account that the federate may encounter problems with the connection to the WSPRC. The WSPRC may also encounter problems with the connection to the federation.

    In case the WSPRC loses the connection to the federation, the Connection Lost \rtiinitiated{} call shall be sent to the Consumer. Any call but the Evoke Multiple Callbacks service shall raise the ``Not Connected'' exception after that.

\end{compactstdpage}

\clearpage

% --- Standard page 332 / source PDF page 346 ---
\begin{compactstdpage}
    \subparagraph{Session timeout}

    Losing the session between the Consumer and the WSPRC corresponds to losing the reference to the RTI ambassador in a federate using the C++ or Java API.

    In case the session between the federate and the WSPRC has become unavailable, all calls to the WSPRC except for connect shall raise the ``Not Connected'' exception.

    In case a Consumers Web Service call to the WSPRC times out, the Consumer may act on its own choosing. It may try to connect later using the session ID already provided. Ultimately it may try to establish a new connection to the same or another WSPRC.

\end{compactstdpage}

